\chapter{Rare Events and Concentration}

\begin{center}\itshape
A rare event is rare in a precise way: its probability has an exponent, and the exponent can be computed.\\
A function of many independent pieces that listens to none of them in particular is almost constant.\\
Between these two facts lies most of modern probability, and a Nobel Prize in physics.
\end{center}

In 2021 the Nobel Prize in Physics went in part to Giorgio Parisi, ``for the discovery of the interplay of disorder and fluctuations in physical systems from atomic to
planetary scales''. Among the discoveries behind the citation is a formula Parisi found in 1979 for the free energy of a toy model of a disordered magnet, the
Sherrington--Kirkpatrick spin glass. The derivation was a feat of physical intuition that no mathematician could follow: it computed the $n$-th power of a random
quantity for integer $n$, then set $n=0$. For a quarter of a century the formula was believed and unproved. Francesco Guerra proved in 2003 that it is an upper bound,
and Michel Talagrand proved in 2006 that it is exact. In 2024 Talagrand received the Abel Prize ``for his groundbreaking contributions to probability theory and
functional analysis, with outstanding applications in mathematical physics and statistics''.

The tools Talagrand forged on the way are the two themes of this chapter. The first is the arithmetic of rare events: Cram\'er's theorem and Sanov's theorem (Tome~10,
Chapter~7) say that the probability of a rare event decays like $e^{-nI}$, and that the exponent $I$ is a variational problem. The second is concentration of measure
(Tome~10, Chapter~8): a Lipschitz function of many independent variables hardly moves. The problems below compute exponents for maxima, positive walks, insurance
companies, sticky Markov chains and a random-energy model of a spin glass, then use bounded differences on colourings, polls, percolation and random formulas, and end in
Gaussian space, where the largest eigenvalue of a random matrix, the maximum of a million Gaussians and the free energy of the Sherrington--Kirkpatrick model at high
temperature all turn out to be almost deterministic. Along the way you will meet a probability of about $10^{-9}$ that one standard inequality bounds by $0.92$, and
another bounds by $7\times10^{-9}$.

\begin{think}
Toss a fair coin a million times, and look for the longest stretch of consecutive tosses in which at least $80\%$ are heads. Is it about $20$ tosses long, about $70$,
or about $700$? Guess before Problem~\ref{pb:windows}.
\end{think}

Throughout, logarithms are natural. For a real variable $X$, $\Lambda(\theta)=\log\mathbb Ee^{\theta X}$ and $\Lambda^*(a)=\sup_\theta(\theta a-\Lambda(\theta))$ are its cumulant
generating function and rate function, as in Tome~10, Chapter~7; for fair coins $\Lambda^*$ is $I(a)=a\log(2a)+(1-a)\log(2(1-a))$ on $[0,1]$. $D(\nu\,\|\,\mu)$ is the
relative entropy, $Z$ (or $g$, $g_{ij}$) a standard Gaussian, $\phi$ its density, $\overline\Phi(x)=P(Z\ge x)$, and $\varepsilon_1,\varepsilon_2,\dots$ are independent fair
signs. We use the Gaussian tail bounds of Tome~10, Chapter~3 (Exercise~3.1): for $x>0$,
\begin{equation}\label{eq:p3gausstail}
\frac x{1+x^2}\,\phi(x)\le\overline\Phi(x)\le\tfrac12e^{-x^2/2}.
\end{equation}

\section{Exact exponents}

Cram\'er's theorem gives the exponent of $P(S_n\ge na)$. The problems of this section ask for the exponent of other events, and each time the answer is a small
variational problem: which is the cheapest way for the rare event to happen?

\begin{exercise}[$\star$]\label{pb:fixedtail}
\emph{Many fixed points.} Let $F_n$ be the number of fixed points of a uniformly random permutation of $\{1,\dots,n\}$.
\begin{enumerate}[label=(\alph*)]
\item Show that $P(F_n\ge k)\le\frac1{k!}$ for $0\le k\le n$.
\item Show that $P(F_n\ge k)\ge\frac1{3\,k!}$ for $0\le k\le n$ with $k\ne n-1$. What happens for $k=n-1$?
\item Deduce that $\log P(F_n\ge k)=-k\log k+k+O(\log k)$ uniformly in $n\ge k+2$, and compare with the tail of the Poisson law of mean $1$. For $n=10$, compute
$P(F_{10}\ge5)$.
\end{enumerate}
\end{exercise}

\begin{exercise}[$\star$]\label{pb:maxexp}
\emph{The maximum costs nothing more.} Let $X_1,X_2,\dots$ be i.i.d.\ copies of a centred variable $X$ satisfying Cram\'er's condition ($\mathbb Ee^{\theta X}<\infty$ for all
$\theta$), $S_k=X_1+\dots+X_k$, and $0<a$ with $P(X>a)>0$.
\begin{enumerate}[label=(\alph*)]
\item Show that $\Lambda^*(a)\le t\,\Lambda^*(a/t)$ for $0<t\le1$, and deduce from the union bound that $P(\max_{k\le n}S_k\ge na)\le n\,e^{-n\Lambda^*(a)}$.
\item Remove the factor $n$: show that $P(\max_{k\le n}S_k\ge na)\le e^{-n\Lambda^*(a)}$.
\item Show that $\frac1n\log P(\max_{k\le n}S_k\ge na)\to-\Lambda^*(a)$.
\end{enumerate}
\end{exercise}

\begin{exercise}[$\star$]\label{pb:positivewalk}
\emph{Staying positive is cheap.} Let $S_k$ be the simple random walk ($\pm1$ steps with probability $\frac12$).
\begin{enumerate}[label=(\alph*)]
\item For $m>0$, show that $P(S_1>0,\dots,S_n>0\mid S_n=m)=\frac mn$.
\item For $0<a<1$, show that $\frac1n\log P(S_1>0,\dots,S_n>0,\ S_n\ge an)\to-I(a)$: at the exponential scale, staying positive is free.
\item Show that $P(S_1>0,\dots,S_{2n}>0)=\frac12\binom{2n}n4^{-n}\sim\frac1{2\sqrt{\pi n}}$: at the polynomial scale, it is not.
\end{enumerate}
\end{exercise}

\begin{exercise}[$\star$]\label{pb:lowentropy}
\emph{A text that compresses.} Roll a fair die with $m$ faces $n$ times, and let $L_n$ be the empirical distribution of the results and $H(L_n)$ its entropy.
\begin{enumerate}[label=(\alph*)]
\item For $0<h<\log m$, show that $\frac1n\log P\bigl(H(L_n)\le h\bigr)\to-(\log m-h)$.
\item Deduce that the number of words of length $n$ over an alphabet of $m$ letters whose empirical entropy is at most $h$ is $e^{nh+o(n)}$, and explain why such words can be
described with about $\frac{nh}{\log2}$ bits.
\end{enumerate}
\end{exercise}

\begin{exercise}[$\star\star$]\label{pb:moderate}
\emph{Between Gauss and Cram\'er.} Let $S_n$ be the simple random walk and $x_n$ a sequence with $\frac{x_n}{\sqrt n}\to\infty$ and $\frac{x_n}n\to0$.
\begin{enumerate}[label=(\alph*)]
\item Show that $P(S_n\ge x_n)\le e^{-x_n^2/2n}$ for every $n$.
\item Show that $\log P(S_n\ge x_n)\sim-\frac{x_n^2}{2n}$.
\item For $n=10^4$ and $x=1000$, compare $-\frac{x^2}{2n}$, $-nI(\frac xn)$, $\log\overline\Phi(\frac x{\sqrt n})$ and the exact value of $\log P(S_n\ge x)$.
\end{enumerate}
\end{exercise}

\begin{exercise}[$\star\star$]\label{pb:windows}
\emph{The Erd\H os--R\'enyi law.} Toss a fair coin $n$ times. For $\frac12<a<1$ let $R_n$ be the largest $\ell$ such that some block of $\ell$ consecutive tosses contains
at least $a\ell$ heads.
\begin{enumerate}[label=(\alph*)]
\item Show that $P\bigl(R_n\ge(1+\varepsilon)\frac{\log n}{I(a)}\bigr)\to0$ for every $\varepsilon>0$.
\item Show that $P\bigl(R_n\le(1-\varepsilon)\frac{\log n}{I(a)}\bigr)\to0$ for every $\varepsilon>0$, so that $\frac{R_n\,I(a)}{\log n}\to1$ in probability.
\item What does this give for $n=10^6$ and $a=0.8$? What does $a=1$ say? (The notebook simulates.)
\end{enumerate}
\end{exercise}

\begin{exercise}[$\star\star$]\label{pb:lundberg}
\emph{Ruin.} An insurance company starts with capital $u>0$. In year $k$ it receives premiums and pays claims, for a net loss $X_k$; the $X_k$ are i.i.d.\ copies of $X$, with
$\mathbb EX<0$, $P(X>0)>0$ and $\mathbb Ee^{\theta X}<\infty$ for all $\theta$. The company is ruined if $S_n=X_1+\dots+X_n>u$ for some $n$; let $\psi(u)$ be the probability of ruin.
\begin{enumerate}[label=(\alph*)]
\item Show that there is a unique $R>0$ with $\mathbb Ee^{RX}=1$ (the \emph{Lundberg exponent}).
\item Show that $\psi(u)\le e^{-Ru}$ for all $u>0$.
\item Show that $\frac1u\log\psi(u)\to-R$: the exponent is exact. Describe the most likely way to be ruined.
\item Compute $\psi(u)$ exactly when $X=\pm1$ with probabilities $p<\frac12$ and $q=1-p$, and $u$ is a positive integer.
\end{enumerate}
\end{exercise}

\begin{exercise}[Journey]\label{pb:rem}
\emph{Derrida's random energy model.} A spin glass with $n$ spins has $2^n$ configurations $\sigma$. In 1980 Bernard Derrida proposed to forget the spins and give each
configuration an independent random energy: let $H_\sigma$, $\sigma\in\{\pm1\}^n$, be independent $N(0,n)$ variables, and for $\beta\ge0$
\[
Z_n(\beta)=\sum_\sigma e^{\beta H_\sigma},\qquad\beta_c=\sqrt{2\log2}.
\]
We show that $\frac1n\log Z_n(\beta)\to f(\beta)$ in probability, where $f(\beta)=\log2+\frac{\beta^2}2$ for $\beta\le\beta_c$ and $f(\beta)=\beta\beta_c$ for $\beta\ge\beta_c$.
\begin{enumerate}[label=(\alph*)]
\item Compute $\mathbb EZ_n(\beta)$ and show that $P\bigl(\frac1n\log Z_n(\beta)\ge\log2+\frac{\beta^2}2+\varepsilon\bigr)\to0$ for every $\varepsilon>0$.
\item Let $M_n=\max_\sigma H_\sigma$. Show that $\frac{M_n}n\to\beta_c$ in probability.
\item For $0<x<\beta_c$, let $N_n(x)$ be the number of $\sigma$ with $H_\sigma\ge nx$. Show that $\frac1n\log\mathbb EN_n(x)\to\log2-\frac{x^2}2$ and that $P\bigl(N_n(x)<\frac12\mathbb EN_n(x)\bigr)\to0$.
Deduce that $P\bigl(\frac1n\log Z_n(\beta)\le f(\beta)-\varepsilon\bigr)\to0$ for every $\beta\ge0$ and $\varepsilon>0$.
\item Let $\beta>\beta_c$ and $\beta_c<a<\beta$. Show that $\mathbb E\bigl[\sum_\sigma e^{\beta H_\sigma}\mathbf 1_{H_\sigma<na}\bigr]\le e^{n(\log2+\beta a-a^2/2)}$, and conclude that
$\frac1n\log Z_n(\beta)\to f(\beta)$ in probability for every $\beta$.
\item Show that $f$ and $f'$ are continuous at $\beta_c$ but $f''$ jumps from $1$ to $0$. Interpret: what does the Gibbs measure $\frac{e^{\beta H_\sigma}}{Z_n}$ look like
on each side of $\beta_c$? (Figure~\ref{fig:rem}.)
\end{enumerate}
\end{exercise}

\begin{exercise}[Journey]\label{pb:markovld}
\emph{Large deviations for a Markov chain.} Let $(X_k)_{k\ge0}$ be a Markov chain on a finite set $E$ with transition matrix $P$ whose entries are all positive, $f:E\to\R$ a
function, and $S_n=f(X_1)+\dots+f(X_n)$. For $\theta\in\R$ let $P_\theta(i,j)=P(i,j)e^{\theta f(j)}$, and let $\lambda(\theta)$ be its Perron root, with positive right and left
eigenvectors $v_\theta$ and $u_\theta$ (Tome~2, Theorem~G.1).
\begin{enumerate}[label=(\alph*)]
\item Show that $\frac1n\log\mathbb E_i\,e^{\theta S_n}\to\Lambda(\theta):=\log\lambda(\theta)$ for every starting state $i$, and that $\Lambda$ is convex.
\item Let $\Lambda^*(a)=\sup_\theta(\theta a-\Lambda(\theta))$. Show that $\limsup\frac1n\log P_i(S_n\ge na)\le-\Lambda^*(a)$ for every $a$ above $\Lambda'(0)$.
\item Show that $Q_\theta(i,j)=\frac{P_\theta(i,j)v_\theta(j)}{\lambda(\theta)v_\theta(i)}$ is a stochastic matrix, with stationary law $\pi_\theta(j)\propto u_\theta(j)v_\theta(j)$, and that
$\Lambda'(\theta)=\sum_j\pi_\theta(j)f(j)$. In particular $\Lambda'(0)$ is the mean of $f$ under the stationary law of $P$.
\item Let $\theta>0$ and $a=\Lambda'(\theta)$. Show that $\liminf\frac1n\log P_i(S_n\ge na-n\delta)\ge-\Lambda^*(a)-\theta\delta$ for every $\delta>0$, and deduce that
$\frac1n\log P_i(S_n\ge na)\to-\Lambda^*(a)$ whenever $\Lambda^*$ is continuous at $a$.
\item The \emph{sticky chain} on $\{0,1\}$ stays put with probability $1-p$ and jumps with probability $p$; take $f(j)=j$, so that $S_n$ is the time spent in state $1$. Show
that $\lambda(\theta)=\frac12\bigl(T+\sqrt{T^2-4(1-2p)e^\theta}\bigr)$ with $T=(1-p)(1+e^\theta)$, and compute $\Lambda^*(0.9)$ for $p=\frac12$, $p=0.1$ and $p=0.02$
(Figure~\ref{fig:markovrate}).
\end{enumerate}
\end{exercise}

\section{Bounded differences and correlations}

McDiarmid's inequality (Tome~10, Theorem~8.3) and the Efron--Stein inequality (Tome~10, Theorem~8.1) only ask how much a function can move when one coordinate changes.
They apply to colourings, to samples, to the shortest path through a random landscape and to random logical formulas. The last problem of the section uses a
different structure, positive correlation, to compute the probability that a random graph has no triangle.

\begin{exercise}[$\star$]\label{pb:discrepancy}
\emph{Colouring to balance.} Let $S_1,\dots,S_n$ be subsets of $\{1,\dots,n\}$. A colouring $\chi\in\{\pm1\}^n$ has \emph{discrepancy} $\max_i|\chi(S_i)|$, where
$\chi(S)=\sum_{j\in S}\chi_j$.
\begin{enumerate}[label=(\alph*)]
\item Show that some colouring has discrepancy at most $\sqrt{2n\log(2n)}$.
\item Let $n$ be a power of $2$ and $H$ the Sylvester--Hadamard matrix of order $n$ ($H_1=(1)$, $H_{2m}=\bigl(\begin{smallmatrix}H_m&H_m\\H_m&-H_m\end{smallmatrix}\bigr)$), and
$S_i=\{j:H_{ij}=1\}$. Show that every colouring has discrepancy at least $\frac{\sqrt n}3$ for these sets.
\end{enumerate}
\end{exercise}

\begin{exercise}[$\star$]\label{pb:blowup}
\emph{Blowing up.} Let $X$ be uniform on the cube $\{0,1\}^n$ and $A$ a non-empty subset, and let $d(x,A)$ be the Hamming distance from $x$ to $A$.
\begin{enumerate}[label=(\alph*)]
\item Show that $\mathbb Ed(X,A)\le\sqrt{\frac n2\log\frac1{P(X\in A)}}$, and that $P\bigl(d(X,A)\ge\mathbb Ed(X,A)+t\bigr)\le e^{-2t^2/n}$.
\item Deduce that if $P(X\in A)\ge\frac12$, then all but a fraction $e^{-2s^2}$ of the cube is within distance $\bigl(s+\sqrt{\log2/2}\bigr)\sqrt n$ of $A$.
\item Show that two subsets of the cube of measure at least $e^{-s}$ each are at Hamming distance less than $\sqrt{2ns}$ from each other.
\end{enumerate}
\end{exercise}

\begin{exercise}[$\star\star$]\label{pb:withoutrepl}
\emph{Polling without replacement.} A population has $N$ members with numbers $c_1,\dots,c_N$. Let $X_1,\dots,X_n$ be a sample \emph{with} replacement ($X_i=c_{I_i}$ with $I_i$
independent and uniform in $\{1,\dots,N\}$), and $Y_1,\dots,Y_n$ a sample \emph{without} replacement ($n\le N$).
\begin{enumerate}[label=(\alph*)]
\item Show that $\mathbb Ef(Y_1+\dots+Y_n)\le\mathbb Ef(X_1+\dots+X_n)$ for every convex function $f:\R\to\R$.
\item Deduce that Hoeffding's inequality holds for $Y_1+\dots+Y_n$: if all $c_j\in[0,1]$ and $\mu$ is their mean, then $P(\sum Y_i-n\mu\le-t)\le e^{-2t^2/n}$.
\item In a town of $10\,000$ voters, $5\,500$ support a project. A pollster asks $1\,000$ distinct voters. Bound the probability that at most half of them support it, and
compare with the exact value.
\end{enumerate}
\end{exercise}

\begin{exercise}[$\star\star$]\label{pb:bernstein}
\emph{Bernstein's inequality.} Let $X_1,\dots,X_n$ be independent with $X_i-\mathbb EX_i\le b$, $S=X_1+\dots+X_n$ and $v=\operatorname{Var}S$.
\begin{enumerate}[label=(\alph*)]
\item Show that $\mathbb Ee^{\theta(X_i-\mathbb EX_i)}\le\exp\bigl(\operatorname{Var}X_i\frac{e^{\theta b}-1-\theta b}{b^2}\bigr)$ for $\theta>0$.
\item Show that $P(S-\mathbb ES\ge t)\le\exp\Bigl(-\dfrac{t^2}{2(v+bt/3)}\Bigr)$ for $t>0$.
\item Let $S$ be binomial with $n=10^6$ and $p=10^{-3}$. Compare Hoeffding's bound, Bernstein's bound and the exact value of $P(S\ge1200)$. Why does Hoeffding fail so
badly?
\end{enumerate}
\end{exercise}

\begin{exercise}[$\star\star$]\label{pb:fpp}
\emph{First passage percolation.} Give each edge $e$ of the grid $\Z^2$ an independent weight $w_e$, uniform on $[1,2]$. The passage time $T(x,y)$ is the smallest total
weight of a path from $x$ to $y$, and $T_n=T\bigl((0,0),(n,0)\bigr)$.
\begin{enumerate}[label=(\alph*)]
\item Show that $T_n$ is a minimum over finitely many paths, that $\mathbb ET_{m+n}\le\mathbb ET_m+\mathbb ET_n$, and that $\frac{\mathbb ET_n}n\to\mu$ for some $1\le\mu\le\frac32$.
\item Show that $\operatorname{Var}T_n\le2n$.
\item Deduce that $\frac{T_n}n\to\mu$ in probability.
\end{enumerate}
\end{exercise}

\begin{exercise}[$\star\star$]\label{pb:ksat}
\emph{Random $k$-SAT.} A random $k$-formula on $n$ Boolean variables has $m$ independent clauses; each clause picks $k$ distinct variables uniformly and negates each
independently with probability $\frac12$, and is the disjunction of the $k$ literals. Let $m=\lfloor\alpha n\rfloor$, and $q=2^{-k}$.
\begin{enumerate}[label=(\alph*)]
\item Show that the formula is satisfiable with probability tending to $0$ if $\alpha>\alpha_1(k)=\frac{\log2}{-\log(1-q)}$. Compute $\alpha_1(3)$, $\alpha_1(4)$ and $\alpha_1(5)$.
\item Let $U$ be the smallest number of clauses violated by an assignment. Show that $P(|U-\mathbb EU|\ge t)\le2e^{-2t^2/m}$.
\item If $\alpha>\alpha_1(k)$, show that there is $\delta>0$ with $P(U\le\delta m)\to0$: the formula is not merely unsatisfiable but far from satisfiable.
\end{enumerate}
\end{exercise}

\begin{exercise}[$\star\star\star$]\label{pb:janson}
\emph{No triangles.} Let $x\in\{0,1\}^N$ have independent coordinates with $P(x_e=1)=p_e$. A function is \emph{increasing} if it does not decrease when a coordinate
goes from $0$ to $1$.
\begin{enumerate}[label=(\alph*)]
\item (Harris.) Show that $\mathbb E[fg]\ge\mathbb Ef\,\mathbb Eg$ for increasing $f,g$, and also for decreasing $f,g$.
\item (Janson.) Let $S_1,\dots,S_m\subseteq\{1,\dots,N\}$, $B_i$ the event that $x_e=1$ for all $e\in S_i$, $\mu=\sum_iP(B_i)$, and
$\Delta=\sum P(B_i\cap B_j)$ over ordered pairs $i\ne j$ with $S_i\cap S_j\ne\emptyset$. Show that
\[
\prod_i\bigl(1-P(B_i)\bigr)\le P\Bigl(\bigcap_i\overline{B_i}\Bigr)\le e^{-\mu+\Delta/2}.
\]
\item In the random graph $G(n,p)$ with $p=\frac cn$, show that $P(\text{no triangle})\to e^{-c^3/6}$.
\item With $p=n^{-3/5}$, show that $\log P(\text{no triangle})\sim-\frac{n^{6/5}}6$.
\end{enumerate}
\end{exercise}

\section{Gaussian worlds}

In Gaussian space, concentration is as clean as it gets: a function with Lipschitz constant $L$ deviates from its mean by $t$ with probability at most $e^{-t^2/2L^2}$, and
its variance is at most $L^2$ (Tome~10, Theorems~8.5 and~8.6). The art is to see the Lipschitz function, and to know when the bound is far from the truth.

\begin{exercise}[$\star$]\label{pb:lambdamax}
\emph{The top eigenvalue.} Let $H$ be a random symmetric $n\times n$ matrix with independent entries on and above the diagonal, $H_{ij}=\frac{g_{ij}}{\sqrt n}$ for $i<j$
and $H_{ii}=\sqrt{\frac2n}\,g_{ii}$ (the Gaussian orthogonal ensemble).
\begin{enumerate}[label=(\alph*)]
\item Show that $\lambda_{\max}(H)$ is a $\sqrt{2/n}$-Lipschitz function of the $\frac{n(n+1)}2$ Gaussians $g_{ij}$.
\item Deduce that $P\bigl(|\lambda_{\max}-\mathbb E\lambda_{\max}|\ge t\bigr)\le2e^{-nt^2/4}$ and $\operatorname{Var}\lambda_{\max}\le\frac2n$.
\item For $n=200$ the notebook finds a standard deviation of about $0.036$. Compare with (b).
\end{enumerate}
\end{exercise}

\begin{exercise}[$\star\star$]\label{pb:sphere}
\emph{L\'evy's sphere.} Let $U$ be uniform on the unit sphere $S^{n-1}\subseteq\R^n$, $n\ge2$.
\begin{enumerate}[label=(\alph*)]
\item Show that every $1$-Lipschitz $f:S^{n-1}\to\R$ (for the Euclidean distance) extends to a $1$-Lipschitz function on $\R^n$.
\item Show that the map $x\mapsto\frac x{\max(|x|,r)}$ is $\frac1r$-Lipschitz on $\R^n$.
\item Deduce that for every $1$-Lipschitz $f$ on $S^{n-1}$ there is a constant $c$ with
\[
P\bigl(|f(U)-c|\ge t\bigr)\le2e^{-nt^2/8}+e^{-(\sqrt{n-1}-\sqrt n/2)^2/2}.
\]
\item Apply this to $f(x)=x_1$, and compare with the exact value of $P(|U_1|\ge t)$ for $n=1000$ and $t=0.1$, $0.2$, $0.3$.
\end{enumerate}
\end{exercise}

\begin{exercise}[$\star\star$]\label{pb:gumbelmax}
\emph{The largest of a million Gaussians.} Let $M_n$ be the maximum of $n$ independent standard Gaussians, $a_n=\sqrt{2\log n}$ and
$b_n=a_n-\frac{\log\log n+\log4\pi}{2a_n}$.
\begin{enumerate}[label=(\alph*)]
\item Show that $\bigl(\frac1x-\frac1{x^3}\bigr)\phi(x)\le\overline\Phi(x)\le\frac{\phi(x)}x$ for $x>0$.
\item Show that $P\bigl(a_n(M_n-b_n)\le x\bigr)\to\exp(-e^{-x})$ for every real $x$ (the Gumbel law).
\item Show that $\operatorname{Var}M_n\le1$. What does (b) suggest instead? For $n=10^6$, compare $a_n$, $b_n$ and $b_n+\frac\gamma{a_n}$ ($\gamma$ is Euler's constant) with
the simulated mean of $M_n$.
\end{enumerate}
\end{exercise}

\begin{exercise}[$\star\star$]\label{pb:achlioptas}
\emph{Projections with coins.} Let $R$ be a $k\times d$ matrix of independent fair signs, $u\in\R^d$ a unit vector, $Q_i=\sum_jR_{ij}u_j$ and $Y=\frac1k|Ru|^2=\frac1k\sum_iQ_i^2$.
\begin{enumerate}[label=(\alph*)]
\item Show that $\mathbb EQ_i^{2m}\le\mathbb EZ^{2m}$ for every $m$, and deduce $\mathbb Ee^{\lambda Q_i^2}\le(1-2\lambda)^{-1/2}$ for $0\le\lambda<\frac12$.
\item Show that $P(Y\ge1+\varepsilon)\le\bigl((1+\varepsilon)e^{-\varepsilon}\bigr)^{k/2}$ and $P(Y\le1-\varepsilon)\le e^{-k\varepsilon^2/6}$ for $0<\varepsilon<1$.
\item Deduce the Johnson--Lindenstrauss lemma with sign matrices: if $k>\frac{24\log N}{\varepsilon^2}$, then for any $N$ points $x_1,\dots,x_N\in\R^d$, with positive probability
$(1-\varepsilon)|x_i-x_j|^2\le\frac1k|Rx_i-Rx_j|^2\le(1+\varepsilon)|x_i-x_j|^2$ for all $i,j$.
\end{enumerate}
\end{exercise}

\begin{exercise}[$\star\star$]\label{pb:salemzygmund}
\emph{Random signs on the circle.} Let $P(t)=\sum_{k=1}^n\varepsilon_ke^{ikt}$ and $\|P\|_\infty=\max_t|P(t)|$.
\begin{enumerate}[label=(\alph*)]
\item Show that $\|P\|_\infty\ge\sqrt n$ for every choice of signs.
\item Show that $\mathbb E\|P\|_\infty\le2\sqrt{n\log(16n^2)}+1$.
\item Show that $P\bigl(|\|P\|_\infty-\mathbb E\|P\|_\infty|\ge s\sqrt n\bigr)\le2e^{-s^2/2}$.
\end{enumerate}
\end{exercise}

\begin{exercise}[Journey/$\star\star\star$]\label{pb:skconc}
\emph{The Sherrington--Kirkpatrick model at high temperature.} Let $g_{ij}$, $1\le i<j\le n$, be independent standard Gaussians, and for $\sigma\in\{\pm1\}^n$ and $\beta>0$
\[
H(\sigma)=\frac\beta{\sqrt n}\sum_{i<j}g_{ij}\sigma_i\sigma_j,\qquad Z_n=\sum_\sigma e^{H(\sigma)},\qquad F_n=\log Z_n .
\]
\begin{enumerate}[label=(\alph*)]
\item Show that $F_n$ is a $\beta\sqrt{\frac{n-1}2}$-Lipschitz function of the $g_{ij}$, and deduce $P(|F_n-\mathbb EF_n|\ge t)\le2e^{-t^2/\beta^2(n-1)}$.
\item Show that $\mathbb EZ_n=2^ne^{\beta^2(n-1)/4}$, so that $\mathbb EF_n\le n\log2+\frac{\beta^2(n-1)}4$.
\item Let $\beta<1$. Show that $\dfrac{\mathbb EZ_n^2}{(\mathbb EZ_n)^2}\le\dfrac{e^{-\beta^2/2}}{\sqrt{1-\beta^2}}$.
\item Deduce that for $\beta<1$, $\frac{F_n}n\to\log2+\frac{\beta^2}4$ almost surely.
\item Show that $\mathbb EF_n\le n\log2+\beta n\sqrt{\log2}$, so that the formula of (d) fails for $\beta>4\sqrt{\log2}$.
\end{enumerate}
\end{exercise}

\section*{Hints}

\begin{hint}{pb:fixedtail}
For (a), $\mathbb E\binom{F_n}k=\frac1{k!}$. For (b), $P(F_n=k)=\frac1{k!}\sum_{j=0}^{n-k}\frac{(-1)^j}{j!}$ (derangements).
\end{hint}

\begin{hint}{pb:maxexp}
For (a), $\Lambda^*$ is convex and $\Lambda^*(0)=0$. For (b), $e^{\theta S_k}$ is a submartingale for $\theta\ge0$; use Doob's maximal inequality.
\end{hint}

\begin{hint}{pb:positivewalk}
(a) is the ballot theorem. For (c), sum (a) over $m$ and compute $\mathbb E|S_{2n}|$ by a telescoping sum.
\end{hint}

\begin{hint}{pb:lowentropy}
$D(\nu\,\|\,\text{uniform})=\log m-H(\nu)$. Use Sanov's theorem.
\end{hint}

\begin{hint}{pb:moderate}
For (b), tilt the coin so that its mean becomes $\frac{x_n+2\sqrt n}n$, and look at the window $[x_n,x_n+4\sqrt n]$.
\end{hint}

\begin{hint}{pb:windows}
For the upper bound, a union bound over starting points and lengths. For the lower bound, cut the $n$ tosses into disjoint blocks and use the lower bound of Tome~3,
Theorem~4.7, in each.
\end{hint}

\begin{hint}{pb:lundberg}
$e^{RS_n}$ is a martingale; stop it at the ruin time. For the lower bound, ruin happens if $S_n>u$ for the single value $n=\lceil tu\rceil$; optimise over $t$.
\end{hint}

\begin{hint}{pb:rem}
In (b) and (c), the variables are independent, so the number of configurations above a level is binomial. In (d), a Gaussian change of variables turns the truncated
mean into a Gaussian tail.
\end{hint}

\begin{hint}{pb:markovld}
Write $\mathbb E_ie^{\theta S_n}=(P_\theta^n\mathbf 1)(i)$ and squeeze $\mathbf 1$ between two multiples of $v_\theta$. For (c), differentiate $u_\theta P_\theta v_\theta=\lambda(\theta)u_\theta v_\theta$. For
(d), change measure from $P$ to $Q_\theta$ and use the ergodic theorem.
\end{hint}

\begin{hint}{pb:discrepancy}
For (a), a random colouring and the union bound. For (b), $HH^T=nI$, so $|H\chi|^2=n^2$; and the first row of $H$ is all ones.
\end{hint}

\begin{hint}{pb:blowup}
$d(\cdot,A)$ changes by at most $1$ when one coordinate changes. Apply McDiarmid's lower tail at the point $d=0$.
\end{hint}

\begin{hint}{pb:withoutrepl}
Generate the sample with replacement from a random permutation $\pi$: the $k$-th distinct value drawn is $c_{\pi(k)}$. Then shuffle the multiplicities, and use Jensen's
inequality for conditional expectations.
\end{hint}

\begin{hint}{pb:bernstein}
$\frac{e^x-1-x}{x^2}=\int_0^1(1-s)e^{sx}\dd s$ is increasing. Then $e^u-1-u\le\frac{u^2}{2(1-u/3)}$ for $0\le u<3$, because $k!\ge2\cdot3^{k-2}$.
\end{hint}

\begin{hint}{pb:fpp}
A path of weight at most $2n$ has at most $2n$ edges. For (b), resampling an edge can lower $T_n$ only if the edge lies on a geodesic, and then by at most $1$.
\end{hint}

\begin{hint}{pb:ksat}
For a fixed assignment, the number of violated clauses is binomial$(m,q)$. For (c), a union bound over the $2^n$ assignments and a Chernoff bound for the lower tail.
\end{hint}

\begin{hint}{pb:janson}
Harris: induction on $N$. Janson: write $P(\bigcap\overline{B_i})=\prod_iP(\overline{B_i}\mid\bigcap_{j<i}\overline{B_j})$, split the conditioning events into those that share a coordinate
with $S_i$ and the others, and use Harris twice.
\end{hint}

\begin{hint}{pb:lambdamax}
Weyl's inequality: $|\lambda_{\max}(A)-\lambda_{\max}(B)|\le\|A-B\|\le\|A-B\|_F$.
\end{hint}

\begin{hint}{pb:sphere}
McShane: $F(x)=\inf_{y\in S^{n-1}}\bigl(f(y)+|x-y|\bigr)$. For (c), if $X$ is a standard Gaussian vector, $\frac X{|X|}$ is uniform on the sphere and $|X|$ is rarely below $\frac{\sqrt n}2$.
\end{hint}

\begin{hint}{pb:gumbelmax}
For (b), $P(M_n\le u)=(1-\overline\Phi(u))^n$; show that $n\overline\Phi(b_n+\frac x{a_n})\to e^{-x}$.
\end{hint}

\begin{hint}{pb:achlioptas}
Expand $\mathbb EQ^{2m}$: each term is $0$ or a product of $u_j^{2r_j}$, and the Gaussian term is at least as large. For the lower tail, $e^{-y}\le1-y+\frac{y^2}2$ for $y\ge0$ and
$\mathbb EQ^4\le3$.
\end{hint}

\begin{hint}{pb:salemzygmund}
For (b), a grid of $4n^2$ points, the bound $|P'|\le\frac{n(n+1)}2$, and the maximum of subgaussian variables. For (c), Tome~10, Exercise~8.5.
\end{hint}

\begin{hint}{pb:skconc}
For (c), $H(\sigma)+H(\tau)$ is Gaussian with a variance depending only on $R=\sum_i\sigma_i\tau_i$; then linearise $e^{sR^2/n}$ with a Gaussian integral. For (d),
Paley--Zygmund and (a).
\end{hint}

\section*{Solutions}

\begin{solution}{pb:fixedtail}
\idea Count $k$-sets of fixed points: on average there are exactly $\frac1{k!}$ of them, whatever $n$. When there is at least one, there is usually exactly one.

\writeup (a) A given set of $k$ points is fixed by $(n-k)!$ permutations, so $\mathbb E\binom{F_n}k=\binom nk\frac{(n-k)!}{n!}=\frac1{k!}$. If $F_n\ge k$ then $\binom{F_n}k\ge1$, so
$P(F_n\ge k)\le\mathbb E\binom{F_n}k=\frac1{k!}$ by Markov's inequality.

(b) Permutations with exactly $k$ fixed points are a choice of the $k$ points and a derangement of the other $n-k$; there are $D_m=m!\sum_{j=0}^m\frac{(-1)^j}{j!}$ derangements
of $m$ points (inclusion--exclusion, Tome~3, Chapter~1). So
\[
P(F_n=k)=\binom nk\frac{D_{n-k}}{n!}=\frac1{k!}\,d_{n-k},\qquad d_m=\sum_{j=0}^m\frac{(-1)^j}{j!}.
\]
Now $d_0=1$, $d_1=0$, and for $m\ge2$ the alternating series with decreasing terms gives $d_m\ge1-1+\frac12-\frac16=\frac13$. So $P(F_n\ge k)\ge P(F_n=k)\ge\frac1{3\,k!}$ if
$n-k\ne1$. For $k=n-1$, $P(F_n=n-1)=0$ (a permutation fixing $n-1$ points fixes all), and $P(F_n\ge n-1)=\frac1{n!}=\frac1{(k+1)!}$: the lower bound fails by a factor $k+1$.

(c) For $n\ge k+2$, (a) and (b) give $-\log3-\log k!\le\log P(F_n\ge k)\le-\log k!$, and Stirling's formula gives $\log k!=k\log k-k+O(\log k)$. For the Poisson law,
$P(\mathrm{Po}(1)\ge k)=e^{-1}\sum_{j\ge k}\frac1{j!}$ lies between $\frac{e^{-1}}{k!}$ and $\frac{e^{-1}}{k!}\sum_{i\ge0}(k+1)^{-i}\le\frac{2e^{-1}}{k!}$: up to constant factors,
fixed points have the Poisson tail \emph{all the way up to $n-2$}, not only for fixed $k$. For $n=10$, $P(F_{10}\ge5)=\sum_{k=5}^{10}\frac{d_{10-k}}{k!}=0.003655$, between
$\frac1{360}$ and $\frac1{120}$.

\tools{Tome~3, Chapters~1 and~2 (derangements, Markov's inequality); Tome~4 (Stirling's formula).}
\leads The rate is superexponential, $\frac{\log P(F_n\ge k)}k\to-\infty$: Cram\'er's theorem has nothing to say, because $F_n$ is not a sum of i.i.d.\ pieces. Tails of the
form $\frac1{k!}$ are the signature of Poisson behaviour, and Chapter~1 of this book showed how close to Poisson $F_n$ really is.
\end{solution}

\begin{solution}{pb:maxexp}
\idea The cheapest way for the walk to reach $na$ before time $n$ is to reach it at time $n$: getting there faster costs more per step. Convexity of $\Lambda^*$ says
exactly this, and Doob's inequality turns it into a bound with no loss at all.

\writeup (a) $\Lambda^*$ is convex, non-negative and $\Lambda^*(0)=0$ (Tome~10, Lemma~7.2, since $\mathbb EX=0$). For $0<t\le1$, $a=t\cdot\frac at+(1-t)\cdot0$, so
$\Lambda^*(a)\le t\Lambda^*(\frac at)+(1-t)\Lambda^*(0)=t\Lambda^*(\frac at)$. By Cram\'er's upper bound (Tome~10, Theorem~7.3(a)), with $t=\frac kn$,
\[
P(S_k\ge na)=P\Bigl(S_k\ge k\frac{na}k\Bigr)\le e^{-k\Lambda^*(na/k)}=e^{-n\,t\Lambda^*(a/t)}\le e^{-n\Lambda^*(a)},
\]
and the union bound over $k\le n$ gives $n\,e^{-n\Lambda^*(a)}$.

(b) For $\theta\ge0$, $\mathbb Ee^{\theta X}\ge e^{\theta\mathbb EX}=1$ by Jensen, so $\mathbb E[e^{\theta S_{k+1}}\mid X_1,\dots,X_k]=e^{\theta S_k}\mathbb Ee^{\theta X}\ge e^{\theta S_k}$: the
process $e^{\theta S_k}$ is a non-negative submartingale. Doob's maximal inequality (Tome~3, Theorem~10.15) gives
\[
P\bigl(\max_{k\le n}S_k\ge na\bigr)=P\bigl(\max_{k\le n}e^{\theta S_k}\ge e^{\theta na}\bigr)\le e^{-\theta na}\mathbb Ee^{\theta S_n}=e^{-n(\theta a-\Lambda(\theta))}.
\]
Since $a>0=\mathbb EX$ and $\Lambda$ is convex with $\Lambda'(0)=0$, $\theta a-\Lambda(\theta)\le0$ for $\theta\le0$, so the supremum over $\theta\ge0$ is $\Lambda^*(a)$.

(c) $\max_{k\le n}S_k\ge S_n$, so the probability is at least $P(S_n\ge na)$, and $\frac1n\log P(S_n\ge na)\to-\Lambda^*(a)$ by Cram\'er's theorem when $0<a<\operatorname{ess\,sup}X$.
(If $a$ is the essential supremum, the same holds with $P(X=a)$, as both sides equal $\log P(X=a)$ in the limit; we leave this boundary case aside.)

\tools{Tome~10, Lemma~7.2 and Theorem~7.3; Tome~3, Theorem~10.15 (Doob).}
\leads The same argument bounds the probability that a random walk with negative drift \emph{ever} exceeds a level: that is Problem~\ref{pb:lundberg}. For Brownian motion the
reflection principle gives the exact answer $P(\max_{s\le1}B_s\ge a)=2P(B_1\ge a)$, a factor $2$ that the exponent does not see.
\end{solution}

\begin{solution}{pb:positivewalk}
\idea Conditioned on its endpoint, the walk is a uniform ballot count, and Bertrand's ballot theorem says it stays positive with probability $\frac mn$, a factor that is
polynomial, hence invisible at the exponential scale.

\writeup (a) Given $S_n=m$, the walk has $a=\frac{n+m}2$ up-steps and $b=\frac{n-m}2$ down-steps in a uniformly random order. By the ballot theorem (Tome~3, Theorem~1.12) it stays
strictly positive with probability $\frac{a-b}{a+b}=\frac mn$.

(b) By (a), $P(S_1>0,\dots,S_n>0,\ S_n\ge an)=\sum_{m\ge an}\frac mnP(S_n=m)$, which lies between $a\,P(S_n\ge an)$ and $P(S_n\ge an)$. Since $\frac1n\log a\to0$, both bounds
have the exponent $-I(a)$ (Tome~3, Theorem~4.7, for $S_n\ge an$, which is a statement about $\frac{n+S_n}2$ heads).

(c) Summing (a) over $m>0$: $P(S_1>0,\dots,S_{2n}>0)=\sum_{m>0}\frac m{2n}P(S_{2n}=m)=\frac{\mathbb E[S_{2n}^+]}{2n}=\frac{\mathbb E|S_{2n}|}{4n}$ by symmetry. Now
$P(S_{2n}=2j)=\binom{2n}{n+j}4^{-n}$, and for $j\ge1$,
\[
\binom{2n-1}{n+j-1}-\binom{2n-1}{n+j}=\binom{2n}{n+j}\Bigl(\frac{n+j}{2n}-\frac{n-j}{2n}\Bigr)=\frac jn\binom{2n}{n+j},
\]
so $\sum_{j\ge1}j\binom{2n}{n+j}=n\binom{2n-1}n=\frac n2\binom{2n}n$ by telescoping. Hence $\mathbb E|S_{2n}|=2\sum_{j\ge1}2j\binom{2n}{n+j}4^{-n}=2n\binom{2n}n4^{-n}$, and
$P(S_1>0,\dots,S_{2n}>0)=\frac12\binom{2n}n4^{-n}\sim\frac1{2\sqrt{\pi n}}$ by Stirling.

\tools{Tome~3, Theorems~1.12 and~4.7; Stirling's formula.}
\leads Staying positive costs a factor $n^{-1/2}$ for every walk with centred steps of finite variance (Sparre Andersen's theorem, 1954, which gives the exact law of the
first time the walk becomes non-positive, whatever the step distribution, as long as it is symmetric and continuous). Polynomial prefactors of this kind are the business of
\emph{sharp} large deviations, of which Problem~\ref{pb:moderate}(c) gives a taste.
\end{solution}

\begin{solution}{pb:lowentropy}
\idea For the uniform law, relative entropy and entropy are the same thing seen from both ends: $D(\nu\,\|\,u)=\log m-H(\nu)$. Sanov's theorem then reads off the answer.

\writeup (a) Let $u$ be the uniform law on $A=\{1,\dots,m\}$. Then $D(\nu\,\|\,u)=\sum_a\nu(a)\log(m\,\nu(a))=\log m-H(\nu)$. Let $\Gamma=\{\nu:H(\nu)\le h\}$. Since $H$ is
continuous on the simplex, $\Gamma$ is closed and contains the open set $\{H<h\}$, so its interior contains $\{H<h\}$. Hence $\inf_{\overline\Gamma}D=\log m-h$, and
$\inf_{\Gamma^\circ}D\le\inf_{H<h}(\log m-H)=\log m-h$, because $H$ takes every value in $[0,\log m]$ (it is continuous on the connected simplex, with $H=0$ at a point mass and
$H=\log m$ at $u$). Sanov's theorem (Tome~10, Theorem~7.5) gives the limit $-(\log m-h)$.

(b) Under the uniform law each word has probability $m^{-n}$, so the number of words with $H(L_n)\le h$ is $m^nP(H(L_n)\le h)=e^{n\log m-n(\log m-h)+o(n)}=e^{nh+o(n)}$. Number
these words; an index takes $\log_2e^{nh+o(n)}=\frac{nh}{\log2}+o(n)$ bits. A text of low empirical entropy is compressible, and one of high entropy is not: almost all texts
have entropy close to $\log m$ and need $n\log_2m$ bits.

\tools{Tome~10, Theorem~7.5 (Sanov); Tome~3, Chapter~4 (entropy).}
\leads This is the method of types behind universal compression: the Lempel--Ziv algorithms of 1977 and 1978 compress any stationary ergodic source to its entropy rate
without knowing the source. This links concentration inequalities with error-correcting codes.
\end{solution}

\begin{solution}{pb:moderate}
\idea Between the central limit theorem ($x_n\approx\sqrt n$) and Cram\'er ($x_n\approx n$) the exponent is the Gaussian one, $\frac{x^2}{2n}$. The upper bound is
Chernoff with the Gaussian choice of $\theta$; the lower bound tilts the coin just enough.

\writeup (a) For $\theta\ge0$, $\mathbb Ee^{\theta S_n}=(\cosh\theta)^n\le e^{n\theta^2/2}$ (compare the series term by term: $(2k)!\ge2^kk!$). Markov's inequality with
$\theta=\frac{x_n}n$ gives $P(S_n\ge x_n)\le e^{-\theta x_n+n\theta^2/2}=e^{-x_n^2/2n}$.

(b) Let $s=\frac{x_n+2\sqrt n}n\to0$ and $\theta=\operatorname{artanh}s$. Under the tilted law $Q$, in which the steps are independent and equal to $+1$ with probability
$\frac{1+s}2$, the density of a path is $\frac{dQ}{dP}=\frac{e^{\theta S_n}}{(\cosh\theta)^n}$. On the window $W=\{x_n\le S_n\le x_n+4\sqrt n\}$,
\[
P(W)=\mathbb E_Q\bigl[e^{-\theta S_n}(\cosh\theta)^n;W\bigr]\ge e^{-\theta(x_n+4\sqrt n)+n\log\cosh\theta}\,Q(W).
\]
Under $Q$, $S_n$ has mean $ns=x_n+2\sqrt n$ and variance $n(1-s^2)\le n$, so $Q(W)\ge1-\frac n{4n}=\frac34$ by Chebyshev. As $s\to0$, $\theta=s+O(s^3)$ and
$\log\cosh\theta=\frac{\theta^2}2+O(\theta^4)$, so
\[
-\theta(x_n+4\sqrt n)+n\log\cosh\theta=-\frac{ns^2}2-2\sqrt n\,s+O(ns^4).
\]
Now $\frac{ns^2}2=\frac{(x_n+2\sqrt n)^2}{2n}\sim\frac{x_n^2}{2n}$, while $\sqrt n\,s\sim\frac{x_n}{\sqrt n}$ and $ns^4\sim\frac{x_n^2}n\cdot\frac{x_n^2}{n^2}$ are both $o\bigl(\frac{x_n^2}n\bigr)$ because
$\frac{x_n}{\sqrt n}\to\infty$ and $\frac{x_n}n\to0$. With (a), $\log P(S_n\ge x_n)\sim-\frac{x_n^2}{2n}$.

(c) The notebook finds $-\frac{x^2}{2n}=-50$, $-nI(0.1)=-50.08$, $\log\overline\Phi(10)=-53.23$ and the exact $\log P(S_{10^4}\ge1000)=-53.21$. The moderate deviation
exponent and Cram\'er's exponent are both right to leading order, and both miss the polynomial factor $\frac1{10\sqrt{2\pi}}$ which the Gaussian tail includes: at this
distance the Gaussian approximation, prefactor and all, is better than either exponent.

\tools{Tome~3, Chapter~4 (Chebyshev, Chernoff); Tome~10, Chapter~7 (tilting).}
\leads The general moderate deviation principle (for i.i.d.\ sums with an exponential moment, and much beyond) is due to Cram\'er himself for the upper part of the range.
Corrections to the Gaussian exponent come from the cumulants, through the \emph{Cram\'er series}: $\log P(S_n\ge x)=-\frac{x^2}{2n}+\frac{x^3}{n^2}\lambda_3+\cdots$ plus the
prefactor, with $\lambda_3=0$ for symmetric steps.
\end{solution}

\begin{solution}{pb:windows}
\idea A block of length $\ell$ is unusually rich in heads with probability $e^{-\ell I(a)}$, and there are about $n$ blocks. The longest rich block is where $ne^{-\ell I(a)}$
crosses $1$: $\ell=\frac{\log n}{I(a)}$.

\writeup (a) Let $\ell_0=\lceil(1+\varepsilon)\frac{\log n}{I(a)}\rceil$. A block of length $\ell$ with at least $a\ell$ heads has probability at most $e^{-\ell I(a)}$ (Tome~3,
Theorem~4.7). There are at most $n$ starting points for each length, so
\[
P(R_n\ge\ell_0)\le\sum_{\ell\ge\ell_0}n\,e^{-\ell I(a)}=\frac{n\,e^{-\ell_0I(a)}}{1-e^{-I(a)}}\le\frac{n^{-\varepsilon}}{1-e^{-I(a)}}\to0.
\]

(b) Let $\ell=\lfloor(1-\varepsilon)\frac{\log n}{I(a)}\rfloor$ and cut the tosses into $\lfloor n/\ell\rfloor$ disjoint blocks of length $\ell$. Each has at least $a\ell$ heads with
probability $p_\ell\ge\frac1{\ell+1}e^{-\ell I(\lceil a\ell\rceil/\ell)}$ (Tome~3, Theorem~4.7). Since $I$ is $C^1$ on $]\frac12,1[$ and $\frac{\lceil a\ell\rceil}\ell\le a+\frac1\ell$, we have
$\ell I(\frac{\lceil a\ell\rceil}\ell)\le\ell I(a)+C$ for large $\ell$, so $p_\ell\ge\frac{e^{-C}}{\ell+1}n^{-(1-\varepsilon)}$. The blocks are independent, so
\[
P(R_n<\ell)\le(1-p_\ell)^{\lfloor n/\ell\rfloor}\le\exp\Bigl(-\Bigl\lfloor\frac n\ell\Bigr\rfloor p_\ell\Bigr)\le\exp\Bigl(-c\frac{n^\varepsilon}{\ell(\ell+1)}\Bigr)\to0,
\]
since $\ell=O(\log n)$.

(c) $I(0.8)=0.8\log1.6+0.2\log0.4=0.19274$, so $\frac{\log10^6}{I(0.8)}=71.7$. The notebook finds $R_{10^6}$ between $55$ and $58$ in three runs: the law of large numbers
holds, but the correction is large (it is of order $\log\log n$, see below). For $a=1$ the problem is the longest run of heads, and $I(1)=\log2$ gives $R_n\approx\log_2n$,
about $20$ for a million tosses (Tome~3, Chapter~4).

\tools{Tome~3, Theorem~4.7 (large deviations for coins, with the lower bound).}
\leads Erd\H os and R\'enyi proved in 1970 that the convergence holds almost surely, and observed that the limit determines the law of the coin: knowing $\frac{R_n}{\log n}$ for
all $a$ recovers $I$, hence the distribution. Deheuvels, Devroye and Lynch (1986) found the next term, of order $\log\log n$. The same statistic, under the name of
\emph{local alignment score}, measures the similarity of DNA sequences, and its Gumbel-type fluctuations are what the BLAST software reports as an E-value.
\end{solution}

\begin{solution}{pb:lundberg}
\idea Tilting by $R$ turns the drift down into a drift up, and $e^{RS_n}$ is the density that does it. Under the tilted law ruin is certain, so the original ruin
probability is just the price of the tilt, $e^{-Ru}$ up to the overshoot.

\writeup (a) $\varphi(\theta)=\mathbb Ee^{\theta X}$ is finite, convex and smooth, with $\varphi(0)=1$ and $\varphi'(0)=\mathbb EX<0$. Since $P(X>\eta)>0$ for some $\eta>0$,
$\varphi(\theta)\ge e^{\theta\eta}P(X>\eta)\to\infty$. So $\varphi$ dips below $1$ and comes back above, and by strict convexity ($X$ is not constant) it crosses $1$ at exactly one
$R>0$.

(b) $M_n=e^{RS_n}$ is a martingale: $\mathbb E[M_{n+1}\mid X_1,\dots,X_n]=M_n\varphi(R)=M_n$. Let $\tau$ be the first $n$ with $S_n>u$. For each $N$, $\tau\wedge N$ is a bounded
stopping time, so $\mathbb EM_{\tau\wedge N}=1$ (Tome~3, Theorem~10.12(a)), and
\[
1\ge\mathbb E[M_\tau;\tau\le N]\ge e^{Ru}P(\tau\le N).
\]
Letting $N\to\infty$, $\psi(u)=P(\tau<\infty)\le e^{-Ru}$.

(c) Write $\Lambda=\log\varphi$, so $\Lambda(R)=0$, and let $c=\Lambda'(R)>0$ ($\Lambda$ is convex, negative on $]0,R[$, zero at $R$, so its slope at $R$ is positive; it is not
$0$ since $\Lambda$ is strictly convex). For $t>0$,
\[
t\Lambda^*\bigl(\tfrac1t\bigr)=\sup_\theta\bigl(\theta-t\Lambda(\theta)\bigr)\ge R-t\Lambda(R)=R,
\]
with equality when $t=t^*=\frac1c$, because then $\theta\mapsto\theta-t^*\Lambda(\theta)$ is concave with derivative $1-\frac{\Lambda'(\theta)}c$, zero at $\theta=R$. Now fix
$\delta>0$ small and let $n=\lceil t^*u\rceil$. Ruin happens if $S_n>u$, and $\frac un\le c$, so $\psi(u)\ge P(\frac{S_n}n>c+\delta)$ for $u$ large enough that... in fact for
all $u$, since $c+\delta>\frac un$. By Cram\'er's lower bound (Tome~10, Theorem~7.3(c)) for the open set $]c+\delta,\infty[$ (assuming $c+\delta$ is below the essential
supremum of $X$, which holds for small $\delta$ since $c=\mathbb E_{\mu_R}X$ is the mean of the tilted law, which is strictly below its essential supremum),
\[
\liminf_{u\to\infty}\frac1u\log\psi(u)\ge\liminf\frac nu\cdot\frac1n\log P\Bigl(\frac{S_n}n>c+\delta\Bigr)\ge-t^*\Lambda^*(c+2\delta),
\]
using that $\Lambda^*$ is increasing beyond the mean. As $\delta\to0$, $t^*\Lambda^*(c+2\delta)\to t^*\Lambda^*(c)=R$ by continuity of $\Lambda^*$ inside its domain. With (b),
$\frac1u\log\psi(u)\to-R$.

The optimisation says how ruin happens: the cumulative loss grows linearly at speed $c=\Lambda'(R)$, which is the mean of the tilted law $e^{Rx-\Lambda(R)}\mu(\dd x)=e^{Rx}\mu(\dd x)$,
and ruin comes at time about $\frac uc$. Bad years do not come in one catastrophe but as a steady run of years that look like a different, losing, business.

(d) Here $\varphi(R)=pe^R+qe^{-R}=1$ gives $e^R=\frac qp$. The walk moves by $\pm1$, so it exceeds $u$ exactly when it reaches $u+1$. By gambler's ruin with an upward
drift (Tome~3, Proposition~10.13, letting the lower barrier go to $-\infty$), the probability of ever reaching level $u+1$ is $(\frac pq)^{u+1}$. So
$\psi(u)=\frac pq\,e^{-Ru}$: Lundberg's bound is off by the constant factor $\frac pq$, the price of the overshoot. With $p=0.4$, $R=\log1.5=0.405$ and $\psi(10)=0.0116$.

\tools{Tome~3, Theorem~10.12 (optional stopping) and Proposition~10.13; Tome~10, Theorem~7.3 (Cram\'er).}
\leads Filip Lundberg introduced the exponent in his Uppsala thesis of 1903, and Cram\'er made the theory rigorous in 1930. The exact asymptotics $\psi(u)\sim Ce^{-Ru}$ come from
renewal theory (the Cram\'er--Lundberg approximation); when claims are heavy-tailed, $R$ does not exist and ruin is caused by one huge claim instead, a dichotomy you met
in Tome~10, Exercise~7.5. Tome~11, Chapter~7 sets the model in continuous time.
\end{solution}

\begin{solution}{pb:rem}
\idea There are $2^n$ energies and each one is above $nx$ with probability $e^{-nx^2/2}$, so about $e^{n(\log2-x^2/2)}$ of them are there, as long as this number is large.
The partition function is dominated by the level $x$ that maximises (number of states) $\times$ (Boltzmann weight), $\log2-\frac{x^2}2+\beta x$. When the best level $x=\beta$
lies beyond the last occupied level $\beta_c$, the system freezes into its few lowest states.

\writeup (a) $\mathbb Ee^{\beta H_\sigma}=e^{n\beta^2/2}$ (Gaussian moment generating function), so $\mathbb EZ_n=2^ne^{n\beta^2/2}$, and by Markov
$P(Z_n\ge e^{n\varepsilon}\mathbb EZ_n)\le e^{-n\varepsilon}\to0$.

(b) Write $H_\sigma=\sqrt n\,g_\sigma$. Upper bound: for $a>\beta_c$, $P(M_n\ge na)\le2^nP(g\ge a\sqrt n)\le2^ne^{-na^2/2}=e^{n(\log2-a^2/2)}\to0$. Lower bound: for $0<a<\beta_c$,
let $q_n=\overline\Phi(a\sqrt n)\ge\frac{a\sqrt n}{1+a^2n}\phi(a\sqrt n)$ by~\eqref{eq:p3gausstail}. By independence,
\[
P(M_n<na)=(1-q_n)^{2^n}\le e^{-2^nq_n},\qquad 2^nq_n\ge\frac{c}{\sqrt n}e^{n(\log2-a^2/2)}\to\infty .
\]
So $\frac{M_n}n\to\beta_c$ in probability.

(c) $N_n(x)$ is binomial with $2^n$ trials and success probability $q_n=\overline\Phi(x\sqrt n)$, so $\mathbb EN_n=2^nq_n$ and, by the bounds~\eqref{eq:p3gausstail},
$\frac1n\log\mathbb EN_n\to\log2-\frac{x^2}2>0$. Its variance is at most its mean, so Chebyshev gives $P(N_n<\frac12\mathbb EN_n)\le\frac4{\mathbb EN_n}\to0$. On the complementary
event,
\[
Z_n(\beta)\ge N_n(x)e^{\beta nx}\ge\tfrac12\mathbb EN_n\,e^{\beta nx},\qquad\text{so}\qquad\frac1n\log Z_n\ge\log2-\frac{x^2}2+\beta x-o(1).
\]
The supremum of $\log2-\frac{x^2}2+\beta x$ over $0<x<\beta_c$ is $\log2+\frac{\beta^2}2$ if $\beta<\beta_c$ (at $x=\beta$) and $\log2-\frac{\beta_c^2}2+\beta\beta_c=\beta\beta_c$ if
$\beta\ge\beta_c$ (as $x\to\beta_c$); in both cases it is $f(\beta)$ (the two formulas agree at $\beta_c$, where both equal $2\log 2$). Choosing $x$ within $\varepsilon/2$ of the
supremum gives the claim.

(d) With $H=\sqrt n\,g$ and the change of variables $y=g-\beta\sqrt n$,
\[
\mathbb E\bigl[e^{\beta H}\mathbf 1_{H<na}\bigr]=\int_{-\infty}^{a\sqrt n}e^{\beta\sqrt n y'}\phi(y')\dd y'=e^{n\beta^2/2}P\bigl(g<(a-\beta)\sqrt n\bigr)=e^{n\beta^2/2}\overline\Phi\bigl((\beta-a)\sqrt n\bigr),
\]
which is at most $e^{n\beta^2/2-n(\beta-a)^2/2}=e^{n(\beta a-a^2/2)}$ since $a<\beta$. Multiply by $2^n$. Now on the event $M_n<na$, whose probability tends to $1$ by (b) since
$a>\beta_c$, $Z_n$ equals the truncated sum, and Markov's inequality shows that the truncated sum exceeds $e^{n(\log2+\beta a-a^2/2+\varepsilon)}$ with probability at most
$e^{-n\varepsilon}$. So with probability tending to $1$, $\frac1n\log Z_n\le\log2+\beta a-\frac{a^2}2+\varepsilon$. As $a\downarrow\beta_c$ the right side tends to
$\log2+\beta\beta_c-\log2+\varepsilon=\beta\beta_c+\varepsilon$. Together with (a) for $\beta\le\beta_c$ and (c) for all $\beta$, $\frac1n\log Z_n\to f(\beta)$ in probability.

(e) $f'(\beta)=\beta$ for $\beta<\beta_c$ and $\beta_c$ for $\beta>\beta_c$: continuous. $f''$ is $1$, then $0$. By (c) and (d), for $\beta<\beta_c$ the Gibbs weight sits on the
$e^{n(\log2-\beta^2/2)}$ configurations with energy near $n\beta$, exponentially many, each of exponentially small weight. For $\beta>\beta_c$ it sits on configurations with
energy within $o(n)$ of the maximum; there are only $e^{o(n)}$ of them, and in fact (beyond this problem) finitely many carry most of the weight. The entropy of the Gibbs
measure, $n(\log 2-\frac{\beta^2}2)$ to leading order, reaches $0$ at $\beta_c$ and stays there: the system is \emph{frozen}.

\tools{Gaussian tail bounds~\eqref{eq:p3gausstail}; Chebyshev's and Markov's inequalities; Tome~10, Chapter~7 (the variational picture).}
\leads Derrida introduced the model in 1980--81 as the limit of spin glasses with $p$-spin interactions as $p\to\infty$. At low temperature, the normalised Gibbs weights of the
top configurations converge to a Poisson--Dirichlet law with parameter $\frac{\beta_c}\beta$, the first instance of the hierarchical structures of Parisi's theory. Problem~\ref{pb:skconc}
starts on the Sherrington--Kirkpatrick model, where the energies are correlated.
\end{solution}

\begin{figure}[ht]
\centering\fig{ch02-rem}
\caption{The random energy model: $\frac1n\log Z_n(\beta)$ for one sample each of $n=12$, $16$ and $20$ spins (dots), and the limit $f(\beta)$ (line), which follows the parabola
$\log2+\frac{\beta^2}2$ (dashed) up to $\beta_c=\sqrt{2\log2}\approx1.18$ and then the tangent line $\beta\beta_c$. Beyond $\beta_c$ the finite systems are still above the limit,
by an amount of order $\frac{\log n}n$.}
\label{fig:rem}
\end{figure}

\begin{solution}{pb:markovld}
\idea For independent steps, the moment generating function of $S_n$ is the $n$-th power of a number. For a Markov chain it is the $n$-th power of a matrix, and the largest
eigenvalue takes the place of the number. Tilting the chain by the Perron eigenvector gives a new Markov chain under which the rare event becomes typical.

\writeup (a) Summing over paths, $\mathbb E_ie^{\theta S_n}=\sum_{x_1,\dots,x_n}\prod_kP(x_{k-1},x_k)e^{\theta f(x_k)}=(P_\theta^n\mathbf 1)(i)$ with $x_0=i$. $P_\theta$ has positive entries, so
by Perron's theorem its spectral radius $\lambda(\theta)>0$ is a simple eigenvalue with a positive eigenvector $v=v_\theta$. Let $\alpha=\min_jv(j)$ and $\beta=\max_jv(j)$. Then
$\frac v\beta\le\mathbf 1\le\frac v\alpha$ entrywise, and $P_\theta^n$ preserves such inequalities (its entries are positive), so $\frac{\lambda^n}\beta v\le P_\theta^n\mathbf 1\le\frac{\lambda^n}\alpha v$.
Hence $\frac1n\log\mathbb E_ie^{\theta S_n}\to\log\lambda(\theta)$. Each $\theta\mapsto\frac1n\log\mathbb E_ie^{\theta S_n}$ is convex (H\"older's inequality), so the limit $\Lambda$ is convex.

(b) For $\theta\ge0$, Markov's inequality gives $P_i(S_n\ge na)\le e^{-n\theta a}\mathbb E_ie^{\theta S_n}$, so $\limsup\frac1n\log P_i(S_n\ge na)\le-(\theta a-\Lambda(\theta))$, and we
take the supremum over $\theta\ge0$. For $\theta<0$, convexity gives $\Lambda(\theta)\ge\Lambda(0)+\theta\Lambda'(0)=\theta\Lambda'(0)$ ($\Lambda$ is differentiable by (c)), so
$\theta a-\Lambda(\theta)\le\theta(a-\Lambda'(0))\le0$ when $a\ge\Lambda'(0)$: negative $\theta$ do not improve the supremum, which is therefore $\Lambda^*(a)$.

(c) $\sum_jQ_\theta(i,j)=\frac{(P_\theta v)(i)}{\lambda v(i)}=1$. With $u=u_\theta$ the positive left eigenvector, $\sum_iu(i)v(i)Q_\theta(i,j)=\frac{(uP_\theta)(j)v(j)}\lambda=u(j)v(j)$,
so $\pi_\theta(j)=\frac{u(j)v(j)}{\sum_ku(k)v(k)}$ is stationary. Since $\lambda(\theta)$ is a simple root of the characteristic polynomial of $P_\theta$, whose coefficients are
smooth in $\theta$, the implicit function theorem makes $\lambda$ smooth, and one can choose $v_\theta$, $u_\theta$ smooth too (as normalised solutions of a linear system of
constant rank). Differentiate $P_\theta v_\theta=\lambda v_\theta$ and multiply on the left by $u_\theta$: the terms with $v_\theta'$ cancel because $u_\theta P_\theta=\lambda u_\theta$,
and since $\frac{\dd}{\dd\theta}P_\theta(i,j)=P_\theta(i,j)f(j)$,
\[
\lambda'\,uv=\sum_{i,j}u(i)P_\theta(i,j)f(j)v(j)=\lambda\sum_ju(j)f(j)v(j),\qquad\text{so}\qquad\Lambda'(\theta)=\frac{\lambda'}\lambda=\sum_j\pi_\theta(j)f(j).
\]
At $\theta=0$, $Q_0=P$ and $\pi_0$ is the stationary law of $P$.

(d) For a path $x_0=i,x_1,\dots,x_n$, the definition of $Q=Q_\theta$ telescopes: $\prod_kQ(x_{k-1},x_k)=\prod_kP(x_{k-1},x_k)\cdot e^{\theta S_n}\lambda^{-n}\frac{v(x_n)}{v(i)}$. So for
any event $A$ depending on the first $n$ steps,
\[
P_i(A)=\mathbb E^Q_i\Bigl[e^{-\theta S_n}\lambda^n\frac{v(i)}{v(X_n)};A\Bigr]\ge\frac\alpha\beta\,\mathbb E^Q_i\bigl[e^{-\theta S_n+n\Lambda(\theta)};A\bigr].
\]
Take $A=\{na-n\delta\le S_n\le na+n\delta\}$; on $A$, $e^{-\theta S_n}\ge e^{-\theta n(a+\delta)}$ (as $\theta>0$). Under $Q$, the chain is irreducible (positive entries), so
$\frac{S_n}n\to\sum_j\pi_\theta(j)f(j)=\Lambda'(\theta)=a$ almost surely by the ergodic theorem (Tome~3, Theorem~8.13), and $Q_i(A)\to1$. Hence
\[
\liminf\frac1n\log P_i(S_n\ge na-n\delta)\ge-\theta a+\Lambda(\theta)-\theta\delta=-\Lambda^*(a)-\theta\delta,
\]
because the concave function $\theta'\mapsto\theta'a-\Lambda(\theta')$ has derivative $0$ at $\theta'=\theta$, so its maximum is $\Lambda^*(a)=\theta a-\Lambda(\theta)$. Replacing $a$ by
$a+\delta$ (that is, using $\theta_\delta$ with $\Lambda'(\theta_\delta)=a+\delta$, which exists by continuity and monotonicity of $\Lambda'$ when $a+\delta$ is in its range), we get
$\liminf\frac1n\log P_i(S_n\ge na)\ge-\Lambda^*(a+\delta)-\theta_\delta\delta$, and letting $\delta\to0$, with (b), the limit is $-\Lambda^*(a)$ where $\Lambda^*$ is continuous.

(e) $P_\theta=\begin{pmatrix}1-p&pe^\theta\\p&(1-p)e^\theta\end{pmatrix}$ has trace $T=(1-p)(1+e^\theta)$ and determinant $(1-p)^2e^\theta-p^2e^\theta=(1-2p)e^\theta$, and its larger
eigenvalue is $\frac12(T+\sqrt{T^2-4\det})$. For $p=\frac12$ the chain is i.i.d., $\lambda=\frac{1+e^\theta}2$, and $\Lambda^*(0.9)=I(0.9)=0.368$. The notebook maximises
$0.9\theta-\Lambda(\theta)$ numerically: $\Lambda^*(0.9)=0.0444$ for $p=0.1$ and $0.0082$ for $p=0.02$. A sticky chain spends $90\%$ of its time in state $1$ at a cost about
$\frac{p}{\text{something}}$ per step, not a fixed cost: it suffices to make the visits to state $0$ short, and the chain is already inclined to stay where it is.

\tools{Tome~2, Theorem~G.1 (Perron); Tome~3, Theorem~8.13 (ergodic theorem); Tome~10, Chapter~7 (tilting, Cram\'er).}
\leads This is the simplest case of the G\"artner--Ellis theorem and of Donsker and Varadhan's large deviations for Markov processes (1975--83). The same Perron root
$\lambda(\theta)$ gives the central limit theorem: $\Lambda''(0)$ is the asymptotic variance of $\frac{S_n}{\sqrt n}$, which for the sticky chain is $\frac{1-p}{4p}$, nine times the
i.i.d.\ value $\frac14$ when $p=0.1$. The free energy of the one-dimensional Ising model is the logarithm of a Perron root in exactly this way (the transfer matrix method).
\end{solution}

\begin{figure}[ht]
\centering\fig{ch02-markovrate}
\caption{The rate function $\Lambda^*(a)$ of the fraction of time spent in state $1$ by the sticky chain, for $p=\frac12$ (independent fair coins), $p=0.1$ and $p=0.02$.
All vanish at $a=\frac12$; the stickier the chain, the flatter the rate function, and the more likely long excursions are. The dots mark $a=0.9$.}
\label{fig:markovrate}
\end{figure}

\begin{solution}{pb:discrepancy}
\idea A random colouring is already good to within a logarithm. A Hadamard matrix shows that the square root cannot be beaten: its rows are orthogonal, so a colouring
cannot be balanced on all of them at once.

\writeup (a) Let $\chi$ be a uniformly random colouring. For each set, $\chi(S_i)$ is a sum of $|S_i|\le n$ independent signs, so Hoeffding's inequality (Tome~3, Theorem~4.9) gives
$P(|\chi(S_i)|\ge t)\le2e^{-t^2/2|S_i|}\le2e^{-t^2/2n}$. With $t=\sqrt{2n\log(2n)}$, the union bound gives $P(\text{some }|\chi(S_i)|\ge t)\le n\cdot2e^{-\log(2n)}=1$; the
inequality is strict because $P(|\chi(S_i)|\ge t)<2e^{-t^2/2n}$ (for instance since $|\chi(S_i)|\le n<t$ for small $n$, or because the Hoeffding bound is never attained). So some
colouring has all $|\chi(S_i)|<t$. (If one wants a clean margin, $t=\sqrt{2n\log(4n)}$ gives probability at most $\frac12$.)

(b) By induction $HH^T=nI$: $H_{2m}H_{2m}^T=\bigl(\begin{smallmatrix}2H_mH_m^T&0\\0&2H_mH_m^T\end{smallmatrix}\bigr)$. So $|H\chi|^2=\chi^THH^T\chi\cdots$ more simply, $|H^T\chi|^2=n|\chi|^2$
and $H^T=H$ for the Sylvester matrix, so $|H\chi|^2=n\cdot n$ and some row has $|(H\chi)_i|\ge\sqrt n$. Now $(H\chi)_i=\chi(S_i)-\chi(S_i^c)=2\chi(S_i)-\chi([n])$, and $S_1=[n]$ because
the first row is all ones. If $|\chi([n])|\ge\frac{\sqrt n}3$, the set $S_1$ already has discrepancy $\ge\frac{\sqrt n}3$. Otherwise
$2|\chi(S_i)|\ge\sqrt n-\frac{\sqrt n}3$, so $|\chi(S_i)|\ge\frac{\sqrt n}3$.

\tools{Tome~3, Theorem~4.9 (Hoeffding); Tome~2 (orthogonal matrices).}
\leads Joel Spencer proved in 1985 that the logarithm is unnecessary: $n$ sets on $n$ points always have a colouring of discrepancy at most $6\sqrt n$ (``six standard deviations
suffice''). His proof was not constructive, and Spencer conjectured that no efficient algorithm could find such a colouring; Nikhil Bansal disproved this in 2010 with a random
walk guided by semidefinite programming, and Lovett and Meka (2012) found a simpler random walk. For $n=64$, (a) gives $24.9$ and (b) gives $2.67$; the notebook tries random and
greedy colourings on the Hadamard sets.
\end{solution}

\begin{solution}{pb:blowup}
\idea Distance to a set is $1$-Lipschitz in each coordinate. McDiarmid's inequality then forbids it to be $0$ with large probability unless its mean is small: a set of
non-negligible measure is close to everything.

\writeup (a) Changing one coordinate of $x$ changes $d(x,A)$ by at most $1$. McDiarmid (Tome~10, Theorem~8.3) with $c_i=1$ gives, for $D=d(X,A)$,
$P(D-\mathbb ED\le-t)\le e^{-2t^2/n}$ and $P(D-\mathbb ED\ge t)\le e^{-2t^2/n}$. With $t=\mathbb ED$ in the first: $P(X\in A)=P(D=0)\le e^{-2(\mathbb ED)^2/n}$, that is
$\mathbb ED\le\sqrt{\frac n2\log\frac1{P(X\in A)}}$.

(b) If $P(X\in A)\ge\frac12$, then $\mathbb ED\le\sqrt{\frac{n\log2}2}$, and $P(D\ge(s+\sqrt{\log2/2})\sqrt n)\le P(D\ge\mathbb ED+s\sqrt n)\le e^{-2s^2}$.

(c) Let $P(A),P(B)\ge e^{-s}$, and $r=\sqrt{2ns}$. By (a), $\mathbb Ed(X,A)\le\sqrt{\frac{ns}2}=\frac r2$, so $P(d(X,A)\ge r)\le P(D\ge\mathbb ED+\frac r2)\le e^{-2(r/2)^2/n}=e^{-s}$, and
the inequality is strict, as in Problem~\ref{pb:discrepancy}, because the exponential bound is not attained for a variable with values in $\{0,\dots,n\}$... to avoid this
subtlety, take $r$ slightly larger: for every $r'>r$, $P(d(X,A)\ge r')<e^{-s}\le P(X\in B)$, so $B$ meets $\{d(\cdot,A)<r'\}$, and $d(A,B)<r'$. Since distances are integers,
$d(A,B)\le\lceil r\rceil$; in the generic case $r\notin\N$, $d(A,B)<r$... We content ourselves with the statement $d(A,B)\le\sqrt{2ns}$ up to rounding, which is the content.

\tools{Tome~10, Theorem~8.3 (McDiarmid).}
\leads This is the \emph{blowing-up lemma} of Ahlswede, G\'acs and K\"orner (1976), proved with a different argument by Margulis, and used in information theory to show that
error probabilities cannot be traded for rate (strong converses). The sharp form is Harper's vertex-isoperimetric inequality on the cube (1966): Hamming balls have the smallest
neighbourhoods.
\end{solution}

\begin{solution}{pb:withoutrepl}
\idea A sample with replacement is a sample without replacement in which some members are counted several times. Shuffling the multiplicities averages them to $1$, and
averaging can only lower a convex function: the sample without replacement is more concentrated.

\writeup (a) Let $I_1,\dots,I_n$ be independent and uniform in $\{1,\dots,N\}$, $D$ the number of distinct values among them, and $R_i\in\{1,\dots,D\}$ the rank of the value $I_i$
in order of first appearance. Let $J_1,\dots,J_D$ be the distinct values in that order. Given the pattern $(R_1,\dots,R_n)$, every injective sequence $(J_1,\dots,J_D)$ is equally
likely, since every vector $(I_1,\dots,I_n)$ has probability $N^{-n}$. Now let $\pi$ be a uniform random permutation of $\{1,\dots,N\}$, independent of the $I_i$. Then
$(\pi(1),\dots,\pi(D))$ is also, given the pattern, a uniform injective sequence, so $(\pi(R_1),\dots,\pi(R_n))$ has the law of $(I_1,\dots,I_n)$ and
\[
X_1+\dots+X_n\overset{\text{law}}=W=\sum_ic_{\pi(R_i)}=\sum_{k=1}^nM_kc_{\pi(k)},\qquad M_k=\#\{i:R_i=k\},
\]
with $M_k=0$ for $k>D$ and $\sum_kM_k=n$. Let $\sigma$ be a uniform permutation of $\{1,\dots,n\}$, independent of everything, and $W'=\sum_{k=1}^nM_{\sigma(k)}c_{\pi(k)}$. Extending
$\sigma$ by the identity on $\{n+1,\dots,N\}$, $W'=\sum_jM_jc_{\pi\sigma^{-1}(j)}$, and $\pi\sigma^{-1}$ is uniform and independent of $(M_j)$; so $W'$ has the law of $W$. But given
$\pi$ and the $M_j$, $\mathbb E[M_{\sigma(k)}]=\frac1n\sum_jM_j=1$, so
\[
\mathbb E[W'\mid\pi,M]=\sum_{k=1}^nc_{\pi(k)},
\]
which is the sum of a sample without replacement. By Jensen's inequality for conditional expectations, $\mathbb Ef(\sum Y_i)=\mathbb Ef(\mathbb E[W'\mid\pi,M])\le\mathbb Ef(W')=\mathbb Ef(\sum X_i)$.

(b) Apply (a) to $f(y)=e^{-\theta y}$, $\theta>0$: $\mathbb Ee^{-\theta\sum Y_i}\le\mathbb Ee^{-\theta\sum X_i}$, and the $X_i$ are independent in $[0,1]$. Hoeffding's proof (Tome~3,
Theorem~4.9) only uses this moment generating function and Markov's inequality, so its bound $e^{-2t^2/n}$ holds for $\sum Y_i$.

(c) With $c_j\in\{0,1\}$, $\mu=0.55$, $n=1000$ and $t=50$: $P(\sum Y_i\le500)\le e^{-2\cdot2500/1000}=e^{-5}=0.0067$. The exact hypergeometric probability is $0.00047$; the bound
is crude (it ignores both the small variance $0.55\times0.45<\frac14$ and the finite population), but it is a guarantee.

\tools{Tome~3, Theorem~4.9 and Chapter~10 (conditional expectation, conditional Jensen).}
\leads Hoeffding proved (a) in his 1963 paper on sums of bounded variables. Serfling (1974) improved the exponent by the factor $1/(1-\frac{n-1}N)$, which matters when the sample
is a large part of the population; it is used for instance to bound errors when a data set is split into training and test sets.
\end{solution}

\begin{solution}{pb:bernstein}
\idea Hoeffding's inequality only knows the range of each variable; Bernstein's also knows the variance. For rare events, like $p=10^{-3}$, the variance is a thousand times
smaller than the square of the range, and the gain is enormous.

\writeup (a) By Taylor's formula with integral remainder, $e^x-1-x=x^2\int_0^1(1-s)e^{sx}\dd s$, so $g(x)=\frac{e^x-1-x}{x^2}$ (with $g(0)=\frac12$) is increasing. Let
$Y=X_i-\mathbb EX_i\le b$ and $\theta>0$. Then $e^{\theta Y}=1+\theta Y+\theta^2Y^2g(\theta Y)\le1+\theta Y+\theta^2Y^2g(\theta b)$, and taking expectations,
$\mathbb Ee^{\theta Y}\le1+\operatorname{Var}X_i\,\theta^2g(\theta b)\le\exp\bigl(\operatorname{Var}X_i\frac{e^{\theta b}-1-\theta b}{b^2}\bigr)$.

(b) For $0\le u<3$, $e^u-1-u=\sum_{k\ge2}\frac{u^k}{k!}\le\frac{u^2}2\sum_{k\ge2}\bigl(\frac u3\bigr)^{k-2}=\frac{u^2}{2(1-u/3)}$, since $k!\ge2\cdot3^{k-2}$. With $u=\theta b$ and (a), by
independence,
\[
\log\mathbb Ee^{\theta(S-\mathbb ES)}\le\frac{v\theta^2}{2(1-b\theta/3)}\qquad(0<\theta<3/b).
\]
Take $\theta=\frac t{v+bt/3}$; then $1-\frac{b\theta}3=\frac v{v+bt/3}$ and Markov's inequality gives
\[
P(S-\mathbb ES\ge t)\le\exp\Bigl(-\frac{t^2}{v+bt/3}+\frac{v\,t^2}{(v+bt/3)^2}\cdot\frac{v+bt/3}{2v}\Bigr)=\exp\Bigl(-\frac{t^2}{2(v+bt/3)}\Bigr).
\]

(c) Here $b=1-p$, $v=np(1-p)=999$ and $t=200$: Bernstein gives $e^{-18.77}=7.1\times10^{-9}$, Hoeffding gives $e^{-2t^2/n}=e^{-0.08}=0.92$, and the exact value is
$4.6\times10^{-10}$. Hoeffding treats each Bernoulli as if it could be anywhere in $[0,1]$, with variance up to $\frac14$, and so it believes that $S$ fluctuates by $\sqrt n/2=500$;
in truth it fluctuates by $\sqrt{999}\approx32$, and $200$ is six of those.

\tools{Taylor's formula; Tome~3, Chapter~4 (Markov, Chernoff).}
\leads Sergei Bernstein published his inequality in the 1920s; Bennett (1962) found the sharper form with $h(u)=(1+u)\log(1+u)-u$, and Freedman (1975) extended it to
martingales, where it is the standard tool for random algorithms and online learning. Bernstein's bound is exactly what the Poisson limit would give: for $t\gg v/b$ the exponent
becomes $\frac{3t}{2b}$, linear rather than quadratic, as for Poisson tails.
\end{solution}

\begin{solution}{pb:fpp}
\idea Subadditivity gives a speed. Efron--Stein bounds the fluctuations by the expected length of the geodesic, because an edge off the geodesic cannot lower the passage
time when it is resampled.

\writeup (a) The straight path has $n$ edges, hence weight at most $2n$; any path of weight at most $2n$ has at most $2n$ edges (each weighs at least $1$), so it stays in the box
$B_n=[-2n,2n]^2$. So $T_n$ is the minimum of the weights of the finitely many self-avoiding paths in $B_n$ with at most $2n$ edges, and it is a function of the finitely many
weights in $B_n$; a minimising path (a \emph{geodesic}) exists, and we fix one, $\gamma_n$, by a deterministic rule (the first in some fixed order of paths). Concatenating paths,
$T((0,0),(m+n,0))\le T((0,0),(m,0))+T((m,0),(m+n,0))$, and the last term has the law of $T_n$ by translation invariance of the weights. So $a_n=\mathbb ET_n$ is subadditive,
and $\frac{a_n}n\to\mu=\inf_n\frac{a_n}n$ by Fekete's lemma (Problem Book~I, Problem~12.21). Every path from $(0,0)$ to $(n,0)$ has at least $n$ edges, so $T_n\ge n$; the
straight path has expected weight $\frac32n$. So $1\le\mu\le\frac32$.

(b) For an edge $e$ of $B_n$, let $T_n^{(e)}$ be the passage time when $w_e$ is replaced by an independent copy $w'_e$. If $e\notin\gamma_n$, then $\gamma_n$ is a path of
the same weight in the new environment, so $T_n^{(e)}\le T_n$. If $e\in\gamma_n$, the new geodesic weighs at least $T_n^{(e)}$ and at most $|w_e-w'_e|\le1$ less in the old
environment, so $T_n\le T_n^{(e)}+1$. Hence $(T_n-T_n^{(e)})_+\le\mathbf 1_{e\in\gamma_n}$. Since $(T_n,T_n^{(e)})$ is exchangeable, $\frac12\mathbb E(T_n-T_n^{(e)})^2=\mathbb E(T_n-T_n^{(e)})_+^2$,
and the Efron--Stein inequality (Tome~10, Theorem~8.1) gives
\[
\operatorname{Var}T_n\le\sum_{e\in B_n}\mathbb E(T_n-T_n^{(e)})_+^2\le\sum_eP(e\in\gamma_n)=\mathbb E|\gamma_n|\le2n .
\]

(c) By Chebyshev, $P(|\frac{T_n}n-\frac{\mathbb ET_n}n|\ge\varepsilon)\le\frac{2}{n\varepsilon^2}\to0$, and $\frac{\mathbb ET_n}n\to\mu$.

\tools{Tome~10, Theorem~8.1 (Efron--Stein); Problem Book~I, Problem~12.21 (Fekete).}
\leads Hammersley and Welsh introduced the model in 1965 and invented subadditive processes for it; Kingman's subadditive ergodic theorem (1968) gives almost sure convergence. The
bound $\operatorname{Var}T_n=O(n)$ is Kesten's (1993); Benjamini, Kalai and Schramm improved it to $O(\frac n{\log n})$ in 2003. The truth is believed to be $n^{2/3}$, the exponent
of the Kardar--Parisi--Zhang universality class, and Tracy--Widom fluctuations; this is proved only for exactly solvable cousins, such as last passage percolation with
exponential weights (Johansson, 2000).
\end{solution}

\begin{solution}{pb:ksat}
\idea Each assignment satisfies each clause with probability $1-q$, so the expected number of satisfying assignments is $2^n(1-q)^m$, which vanishes beyond $\alpha_1(k)$. The
minimum number of violated clauses moves by at most $1$ when a clause changes, so it is concentrated; and beyond $\alpha_1$ it is proportional to $m$.

\writeup (a) A clause is violated by a fixed assignment exactly when its $k$ literals are all false, which happens with probability $2^{-k}=q$ (each negation is a fair coin).
Clauses are independent, so the assignment satisfies all of them with probability $(1-q)^m$, and the expected number of satisfying assignments is
$2^n(1-q)^{\lfloor\alpha n\rfloor}=e^{n(\log2+\alpha\log(1-q))+O(1)}\to0$ when $\alpha>\alpha_1(k)$. By Markov's inequality the formula is unsatisfiable with probability tending
to $1$. Numerically $\alpha_1(3)=5.19$, $\alpha_1(4)=10.74$, $\alpha_1(5)=21.83$; since $-\log(1-q)>q$, $\alpha_1(k)<2^k\log2$, and the two agree to within $\frac12\log2$.

(b) $U$ is a function of the $m$ independent clauses. Replacing one clause changes the number of violated clauses of every assignment by at most $1$, hence their minimum by
at most $1$. McDiarmid with $c_i=1$ gives $P(|U-\mathbb EU|\ge t)\le2e^{-2t^2/m}$.

(c) For a fixed assignment, the number of violated clauses is binomial$(m,q)$, and for $\delta<q$, $P(\text{Bin}(m,q)\le\delta m)\le e^{-mD(\delta\|q)}$ with
$D(\delta\|q)=\delta\log\frac\delta q+(1-\delta)\log\frac{1-\delta}{1-q}$ (Tome~3, Theorem~4.7, by the symmetry stated there). By the union bound over $2^n$ assignments,
\[
P(U\le\delta m)\le2^ne^{-mD(\delta\|q)}=e^{n(\log2-\alpha D(\delta\|q))+O(1)}.
\]
At $\delta=0$, $D(0\|q)=-\log(1-q)$ and $\alpha D(0\|q)>\log2$; by continuity there is $\delta>0$ with $\alpha D(\delta\|q)>\log2$, and then $P(U\le\delta m)\to0$. With (b), $U$ is
in fact within $O(\sqrt m)$ of its mean, which is at least $\delta m-O(\sqrt m)$.

\tools{Tome~3, Theorem~4.7; Tome~10, Theorem~8.3 (McDiarmid).}
\leads The first moment bound is not the truth: for $k=3$ the satisfiability threshold is believed, from the cavity method of statistical physics (M\'ezard, Parisi and Zecchina,
2002), to be $\alpha_3\approx4.267$, below $5.19$. Friedgut (1999) proved that the transition is sharp, Achlioptas and Peres (2004) that the threshold is $2^k\log2-O(k)$, and Ding,
Sly and Sun (2022) proved the physicists' exact formula for all large $k$. The case $k=3$ is open.
\end{solution}

\begin{solution}{pb:janson}
\idea Harris: increasing events help each other. Janson: avoiding all the $B_i$ is almost as likely as if they were independent, up to a correction $\Delta$ that only counts
overlapping pairs. When the pairs overlap little, the probability of seeing none is $e^{-\mu}$, as for a Poisson count.

\writeup (a) For $N=1$: $f$ and $g$ are increasing functions of one bit, so $(f(1)-f(0))(g(1)-g(0))\ge0$, and
$\mathbb Efg-\mathbb Ef\,\mathbb Eg=p(1-p)(f(1)-f(0))(g(1)-g(0))\ge0$ (expand both sides). For $N>1$, condition on $x_N$: for each value of $x_N$, $f$ and $g$ are increasing
functions of $x_1,\dots,x_{N-1}$, which are independent of $x_N$, so by induction $\mathbb E[fg\mid x_N]\ge\mathbb E[f\mid x_N]\,\mathbb E[g\mid x_N]$. The two conditional expectations are
increasing functions of $x_N$, so the case $N=1$ gives $\mathbb E\bigl[\mathbb E[f\mid x_N]\mathbb E[g\mid x_N]\bigr]\ge\mathbb Ef\,\mathbb Eg$. For decreasing $f,g$ apply this to $-f,-g$. In particular,
an increasing and a decreasing event are negatively correlated ($\mathbf 1_A$ and $-\mathbf 1_B$ are both increasing).

(b) Lower bound: $\overline{B_i}$ are decreasing events, so by (a) and induction, $P(\bigcap\overline{B_i})\ge\prod P(\overline{B_i})$. Upper bound: write
$P(\bigcap_i\overline{B_i})=\prod_i\bigl(1-P(B_i\mid C_i)\bigr)$ with $C_i=\bigcap_{j<i}\overline{B_j}$. Fix $i$, and split $C_i=A\cap C$, where $A$ is the intersection over $j<i$
with $S_j\cap S_i\ne\emptyset$, and $C$ over the other $j<i$. Then
\[
P(B_i\mid A\cap C)\ge\frac{P(B_i\cap A\cap C)}{P(C)}=P(B_i\cap A\mid C)\ge P(B_i\mid C)-\sum_{j<i,\,j\sim i}P(B_i\cap B_j\mid C),
\]
where $j\sim i$ means $S_j\cap S_i\ne\emptyset$. Now $B_i$ depends only on the coordinates in $S_i$ and $C$ only on others, so $P(B_i\mid C)=P(B_i)$. And $B_i\cap B_j$ is increasing
while $C$ is decreasing, so $P(B_i\cap B_j\cap C)\le P(B_i\cap B_j)P(C)$ by (a), that is, $P(B_i\cap B_j\mid C)\le P(B_i\cap B_j)$. Hence, using $1-x\le e^{-x}$,
\[
P\Bigl(\bigcap_i\overline{B_i}\Bigr)\le\prod_i\Bigl(1-P(B_i)+\sum_{j<i,\,j\sim i}P(B_i\cap B_j)\Bigr)\le\exp\Bigl(-\mu+\sum_i\sum_{j<i,\,j\sim i}P(B_i\cap B_j)\Bigr)=e^{-\mu+\Delta/2}.
\]

(c) The coordinates are the $\binom n2$ edges, $S_i$ the three edges of the $i$-th triangle. Then $\mu=\binom n3p^3\to\frac{c^3}6$. Two distinct triangles sharing an edge span $4$
vertices and $5$ edges; there are at most $\binom n3\cdot3(n-3)\le\frac{n^4}2$ ordered pairs, so $\Delta\le\frac{n^4}2p^5=\frac{c^5}{2n}\to0$ (triangles sharing one vertex share no
edge and do not count). So $P(\text{no triangle})\le e^{-\mu+\Delta/2}\to e^{-c^3/6}$. The lower bound $(1-p^3)^{\binom n3}\to e^{-c^3/6}$ as well.

(d) Now $\mu=\binom n3p^3\sim\frac{n^3}6n^{-9/5}=\frac{n^{6/5}}6$ and $\Delta\le\frac{n^4}2n^{-3}=\frac n2=o(\mu)$, so $\log P\le-\mu(1+o(1))$. The Harris bound gives
$\log P\ge\binom n3\log(1-p^3)=-\binom n3p^3(1+O(p^3))=-\mu(1+o(1))$. So $\log P(\text{no triangle})\sim-\frac{n^{6/5}}6$, although the number of triangles is now huge and
nothing is Poisson any more.

\tools{Induction, conditional expectations (Tome~3, Chapter~10); Problem Book~I, Problem~8.12 (the uniform case of Harris's inequality).}
\leads Harris proved his inequality in 1960 for percolation, and Kleitman independently in 1966 for set systems. Janson's inequality is from 1990 (with \L uczak and Ruci\'nski);
the short proof above is Boppana and Spencer's (1989). When $\Delta\ge\mu$ Janson also proved $P(\bigcap\overline{B_i})\le e^{-\mu^2/2\Delta}$, which gives the exponent for every
$p$ up to constants. The opposite question, the probability of \emph{too many} triangles, is much harder (the ``infamous upper tail''); its exponent was determined only in the
2010s, by Chatterjee, Dembo, Harel, Mousset, Samotij and others.
\end{solution}

\begin{solution}{pb:lambdamax}
\idea Eigenvalues of symmetric matrices are $1$-Lipschitz for the operator norm, hence for the Frobenius norm, and the Frobenius norm of $H$ is a scaled Euclidean norm of
the Gaussians. Gaussian concentration does the rest.

\writeup (a) Let $H,H'$ come from Gaussians $g,g'$. Then
\[
\|H-H'\|_F^2=\sum_i\frac2n(g_{ii}-g'_{ii})^2+2\sum_{i<j}\frac1n(g_{ij}-g'_{ij})^2=\frac2n|g-g'|^2 .
\]
By Weyl's inequality (Tome~2, Corollary~E.7), $|\lambda_{\max}(H)-\lambda_{\max}(H')|\le\|H-H'\|\le\|H-H'\|_F$, since the operator norm of a symmetric matrix is its largest
$|\text{eigenvalue}|$, at most the square root of the sum of the squared eigenvalues, which is $\|\cdot\|_F$. So $\lambda_{\max}$ is $\sqrt{2/n}$-Lipschitz in $g$.

(b) Gaussian concentration (Tome~10, Theorem~8.6) with $L^2=\frac2n$ gives $P(\lambda_{\max}-\mathbb E\lambda_{\max}\ge t)\le e^{-t^2/2L^2}=e^{-nt^2/4}$, the same for the lower
tail, and $\operatorname{Var}\lambda_{\max}\le L^2=\frac2n$.

(c) The bound says the standard deviation is at most $\sqrt{2/200}=0.1$; the truth, $0.036$, is about three times smaller. The gap grows with $n$: the true fluctuations are of
order $n^{-2/3}$, not $n^{-1/2}$.

\tools{Tome~2, Corollary~E.7 (Weyl); Tome~10, Theorem~8.6 (Gaussian concentration).}
\leads The mean $\mathbb E\lambda_{\max}$ tends to $2$, the edge of Wigner's semicircle (Tome~11, Chapter~10), and $n^{2/3}(\lambda_{\max}-2)$ converges to the Tracy--Widom law
(1996 for this ensemble), which also governs the longest increasing subsequence and the KPZ growth models of Problem~\ref{pb:fpp}. Concentration inequalities that are off by a
power of $n$ like this one are the subject of \emph{superconcentration} (Chatterjee, 2014): the maximum of many weakly correlated variables fluctuates much less than any
single one of them.
\end{solution}

\begin{solution}{pb:sphere}
\idea A uniform point on the sphere is a Gaussian vector divided by its length, and the length is almost $\sqrt n$. Extend $f$ from the sphere to space, compose with the radial
projection, and Gaussian concentration transfers to the sphere with Lipschitz constant $\frac1{\sqrt n}$, up to the factor lost in the extension.

\writeup (a) McShane's extension: $F(x)=\inf_{y\in S^{n-1}}(f(y)+|x-y|)$. For $x\in S^{n-1}$, $F(x)\le f(x)$ (take $y=x$) and $f(y)+|x-y|\ge f(x)$ since $f$ is $1$-Lipschitz, so $F=f$ on the
sphere. For any $x,x'$ and $y$, $f(y)+|x-y|\le f(y)+|x'-y|+|x-x'|$; taking infima, $F(x)\le F(x')+|x-x'|$, and by symmetry $F$ is $1$-Lipschitz.

(b) The map $p(x)=x$ for $|x|\le r$ and $p(x)=\frac{rx}{|x|}$ for $|x|>r$ is the nearest-point projection onto the closed ball of radius $r$, and the projection onto a closed
convex set is $1$-Lipschitz (for $a=p(x)$, $b=p(y)$: $\langle x-a,b-a\rangle\le0$ and $\langle y-b,a-b\rangle\le0$; adding, $|a-b|^2\le\langle x-y,a-b\rangle\le|x-y||a-b|$). Our map
is $\frac1rp$.

(c) Let $X$ be a standard Gaussian vector in $\R^n$; $\frac X{|X|}$ is uniform on the sphere (the law of $X$ is rotation invariant and so is that of $\frac X{|X|}$, and the uniform
law is the only rotation-invariant law on the sphere). Let $r=\frac{\sqrt n}2$ and $G(x)=F\bigl(\frac x{\max(|x|,r)}\bigr)$, which is $\frac1r$-Lipschitz by (a) and (b), and $c=\mathbb EG(X)$.
On $\{|X|\ge r\}$, $G(X)=f(\frac X{|X|})$. So
\[
P(|f(U)-c|\ge t)\le P(|G(X)-c|\ge t)+P(|X|<r)\le2e^{-r^2t^2/2}+P(|X|<r).
\]
The first term is $2e^{-nt^2/8}$ by Gaussian concentration. For the second, $|x|$ is $1$-Lipschitz, so $\operatorname{Var}|X|\le1$ and $(\mathbb E|X|)^2\ge\mathbb E|X|^2-1=n-1$; then
$P(|X|<r)\le P(|X|-\mathbb E|X|\le-(\sqrt{n-1}-r))\le e^{-(\sqrt{n-1}-\sqrt n/2)^2/2}$ for $n\ge2$.

(d) For $f(x)=x_1$ take $F(x)=x_1$ directly; then $G$ is odd, $c=0$, and $P(|U_1|\ge t)\le2e^{-nt^2/8}+e^{-(\sqrt{n-1}-\sqrt n/2)^2/2}$. Exactly, $U_1^2$ has the Beta law with
parameters $\frac12$ and $\frac{n-1}2$. For $n=1000$ the notebook finds $P(|U_1|\ge0.1)=0.0015$ (bound $0.57$), $P(|U_1|\ge0.2)=1.7\times10^{-10}$ (bound $0.013$) and
$P(|U_1|\ge0.3)=2.9\times10^{-22}$ (bound $2.6\times10^{-5}$). The bound has the right form but the wrong constant: the truth is close to $e^{-nt^2/2}$, and the factor $4$ in the
exponent was lost by the crude choice $r=\frac{\sqrt n}2$ instead of $r\approx\sqrt n$.

\tools{Tome~10, Theorem~8.6 and Exercise~8.3; Tome~8 or Tome~6 (rotation invariance).}
\leads Paul L\'evy proved in 1919--22 that among subsets of the sphere of given area, caps have the smallest $\varepsilon$-neighbourhoods, which gives the sharp constant. Vitali
Milman's use of it in 1971 for Dvoretzky's theorem made concentration of measure a central idea of high-dimensional geometry; Tome~10, Chapter~8 tells the story.
\end{solution}

\begin{solution}{pb:gumbelmax}
\idea The maximum of $n$ Gaussians sits where $n\overline\Phi(u)\approx1$, and Mills' ratio says how $\overline\Phi$ behaves there: shifting $u$ by $\frac x{a_n}$ multiplies
$\overline\Phi(u)$ by $e^{-x}$. That gives the Gumbel law, with fluctuations of order $\frac1{\sqrt{\log n}}$, far smaller than the bound $1$ of Gaussian concentration.

\writeup (a) For $t\ge x>0$, $\frac tx\ge1$, so $\overline\Phi(x)\le\int_x^\infty\frac tx\phi(t)\dd t=\frac{\phi(x)}x$. For the lower bound,
$\frac{\dd}{\dd t}\bigl[(\frac1t-\frac1{t^3})\phi(t)\bigr]=\bigl(-\frac1{t^2}+\frac3{t^4}-1+\frac1{t^2}\bigr)\phi(t)=-\bigl(1-\frac3{t^4}\bigr)\phi(t)\ge-\phi(t)$, and integrating from $x$
to $\infty$ gives $(\frac1x-\frac1{x^3})\phi(x)\le\overline\Phi(x)$.

(b) Let $u_n=b_n+\frac x{a_n}$. Then $P(a_n(M_n-b_n)\le x)=(1-\overline\Phi(u_n))^n$, and it suffices to show $n\overline\Phi(u_n)\to e^{-x}$, because $\overline\Phi(u_n)\to0$ and
$n\log(1-y_n)=-ny_n(1+o(1))$ for $y_n\to0$. By (a), $\overline\Phi(u_n)\sim\frac{\phi(u_n)}{u_n}$ and $u_n\sim a_n$. Now, with $\ell=\log\log n+\log4\pi$,
\[
u_n^2=b_n^2+\frac{2b_nx}{a_n}+\frac{x^2}{a_n^2}=a_n^2-\ell+\frac{\ell^2}{4a_n^2}+2x-\frac{\ell x}{a_n^2}+\frac{x^2}{a_n^2}=2\log n-\ell+2x+o(1).
\]
So
\[
n\frac{\phi(u_n)}{u_n}=\frac{n\,e^{-\log n+\ell/2-x+o(1)}}{\sqrt{2\pi}\,a_n(1+o(1))}=\frac{\sqrt{4\pi\log n}\,e^{-x}}{\sqrt{2\pi}\sqrt{2\log n}}(1+o(1))\to e^{-x}.
\]

(c) $\max$ is $1$-Lipschitz (Tome~10, Example~8.8), so $\operatorname{Var}M_n\le1$ by the Gaussian Poincar\'e inequality. But (b) says $M_n=b_n+\frac{G}{a_n}+o_P(\frac1{a_n})$ with $G$
Gumbel, whose variance is $\frac{\pi^2}6$: if the variances converge too (they do, by a uniform integrability argument we do not give), $\operatorname{Var}M_n\sim\frac{\pi^2}{12\log n}$. For
$n=10^6$: $a_n=5.257$, $b_n=4.766$, $b_n+\frac\gamma{a_n}=4.876$ (the Gumbel law has mean $\gamma$), and the notebook's simulated mean is $4.877$, with standard deviation $0.255$
against $\frac\pi{\sqrt6\,a_n}=0.244$. The first-order answer $\sqrt{2\log n}$ of Tome~10, Exercise~8.4, is off by $0.38$, which is one and a half standard deviations.

\tools{Tome~10, Theorem~8.5 and Example~8.8.}
\leads Fisher and Tippett (1928) found the three possible limit laws of maxima, Gnedenko (1943) their domains of attraction; the Gumbel law is the one for light tails. The slow
convergence ($b_n$ differs from $a_n$ by a $\log\log n$ term) is why extreme-value statistics fits the three-parameter family rather than using the limit directly. Chatterjee's
superconcentration theory explains the gap between $1$ and $\frac{\pi^2}{12\log n}$ through the chaos of the maximiser under small perturbations.
\end{solution}

\begin{solution}{pb:achlioptas}
\idea Even moments of a sum of signs are at most those of a Gaussian with the same variance, so a sign matrix projects at least as well as a Gaussian one. The lower tail
needs only the fourth moment.

\writeup (a) $Q=\sum_jR_ju_j$ and $Z=\sum_jg_ju_j$ (a standard Gaussian, since $|u|=1$). Expanding, $\mathbb EQ^{2m}=\sum\binom{2m}{a}\prod_ju_j^{a_j}\mathbb E\prod_jR_j^{a_j}$, the
sum over multi-indices with $\sum a_j=2m$. The expectation is $1$ if all $a_j$ are even and $0$ otherwise. The same expansion for $Z$ has $\mathbb E\prod g_j^{a_j}=\prod(a_j-1)!!\ge1$
when all $a_j$ are even, $0$ otherwise, and in the even case $\prod u_j^{a_j}\ge0$. So $\mathbb EQ^{2m}\le\mathbb EZ^{2m}$. Summing the exponential series (monotone convergence),
$\mathbb Ee^{\lambda Q^2}=\sum_m\frac{\lambda^m\mathbb EQ^{2m}}{m!}\le\mathbb Ee^{\lambda Z^2}=(1-2\lambda)^{-1/2}$.

(b) Upper tail: by independence of the rows and Markov, for $0<\lambda<\frac12$, $P(kY\ge k(1+\varepsilon))\le e^{-\lambda k(1+\varepsilon)}(1-2\lambda)^{-k/2}$; this is the bound of
the Gaussian case, and $\lambda=\frac\varepsilon{2(1+\varepsilon)}$ gives $((1+\varepsilon)e^{-\varepsilon})^{k/2}$ as in Tome~10, Lemma~8.9. Lower tail: for $y\ge0$, $e^{-y}\le1-y+\frac{y^2}2$,
and $\mathbb EQ^2=1$, $\mathbb EQ^4\le\mathbb EZ^4=3$, so $\mathbb Ee^{-\lambda Q^2}\le1-\lambda+\frac32\lambda^2\le e^{-\lambda+3\lambda^2/2}$. Then
\[
P(kY\le k(1-\varepsilon))\le e^{\lambda k(1-\varepsilon)}\bigl(\mathbb Ee^{-\lambda Q^2}\bigr)^k\le e^{k(-\lambda\varepsilon+3\lambda^2/2)}=e^{-k\varepsilon^2/6}\quad\text{for }\lambda=\frac\varepsilon3 .
\]

(c) Since $\log(1+\varepsilon)-\varepsilon\le-\frac{\varepsilon^2}2+\frac{\varepsilon^3}3$, the upper tail is at most $e^{-k(\varepsilon^2/4-\varepsilon^3/6)}\le e^{-k\varepsilon^2/12}$ for $\varepsilon\le1$, and so is
the lower tail. Apply (b) to $u=\frac{x_i-x_j}{|x_i-x_j|}$ for the fewer than $\frac{N^2}2$ pairs: the probability that some pair is distorted is less than
$N^2e^{-k\varepsilon^2/12}<1$ when $k>\frac{24\log N}{\varepsilon^2}$.

\tools{Tome~10, Lemma~8.9 and Theorem~8.10.}
\leads Dimitris Achlioptas proved this in 2001--03 (``database-friendly random projections''), together with a sparse version in which two thirds of the entries are $0$, which
speeds up the projection threefold. Kane and Nelson (2014) pushed sparsity to $O(\varepsilon^{-1}\log N)$ non-zero entries per column. Larsen and Nelson (2017) proved that
$k=O(\varepsilon^{-2}\log N)$ cannot be improved for any linear or non-linear map.
\end{solution}

\begin{solution}{pb:salemzygmund}
\idea On a fine grid the values of $P$ are subgaussian with variance $n$, and the maximum of $M$ subgaussians is at most $\sqrt{2n\log M}$. The grid needs only polynomially
many points, because $P$ cannot move fast. The lower bound is the energy: the average of $|P|^2$ is $n$.

\writeup (a) By Parseval, $\frac1{2\pi}\int_0^{2\pi}|P(t)|^2\dd t=\sum_k\varepsilon_k^2=n$, so $\max|P|^2\ge n$.

(b) Let $N=4n^2$ and $t_j=\frac{2\pi j}N$. Since $|P'(t)|=|\sum k\varepsilon_ke^{ikt}|\le\frac{n(n+1)}2$, every $t$ is within $\frac\pi N$ of some $t_j$, and
$|P(t)|\le|P(t_j)|+\frac{n(n+1)}2\cdot\frac\pi{4n^2}\le|P(t_j)|+\frac\pi4$. For each $j$, $\operatorname{Re}P(t_j)=\sum_k\varepsilon_k\cos(kt_j)$ satisfies
$\mathbb Ee^{\lambda\operatorname{Re}P(t_j)}=\prod_k\cosh(\lambda\cos kt_j)\le e^{\lambda^2n/2}$, and similarly for $-\operatorname{Re}P$, $\pm\operatorname{Im}P$. For $M$ variables $V_l$ with
$\mathbb Ee^{\lambda V_l}\le e^{\lambda^2n/2}$, Jensen gives $e^{\lambda\mathbb E\max V_l}\le\mathbb Ee^{\lambda\max V_l}\le\sum_l\mathbb Ee^{\lambda V_l}\le Me^{\lambda^2n/2}$, so
$\mathbb E\max V_l\le\frac{\log M}\lambda+\frac{\lambda n}2=\sqrt{2n\log M}$ for the best $\lambda$. With the $M=4N=16n^2$ variables $\pm\operatorname{Re}P(t_j)$, $\pm\operatorname{Im}P(t_j)$, and
$|z|\le\sqrt2\max(|\operatorname{Re}z|,|\operatorname{Im}z|)$,
\[
\mathbb E\|P\|_\infty\le\frac\pi4+\sqrt2\sqrt{2n\log(16n^2)}\le2\sqrt{n\log(16n^2)}+1 .
\]

(c) $P=\sum_k\varepsilon_kv_k$ with $v_k(t)=e^{ikt}$ in the space of continuous functions with the sup norm, $\|v_k\|=1$. Tome~10, Exercise~8.5, gives
$P(|\|P\|_\infty-\mathbb E\|P\|_\infty|\ge s\sqrt n)\le2e^{-s^2n/2n}=2e^{-s^2/2}$.

For $n=1000$ the notebook finds a mean of $89$ for $\|P\|_\infty$, against $\sqrt{n\log n}=83$, and the bound (b) of $258$: the logarithm is right, the constant is not.

\tools{Parseval (Tome~4, Chapter~6); Tome~10, Exercise~8.5.}
\leads Salem and Zygmund proved in 1954 that $\|P\|_\infty$ is of order $\sqrt{n\log n}$ almost surely, and Hal\'asz (1973) that $\frac{\|P\|_\infty}{\sqrt{n\log n}}\to1$. Signs can be
chosen much better than at random: the Rudin--Shapiro polynomials of Problem Book~I, Problem~1.21, have $\|P\|_\infty\le\sqrt{2n}$, and Balister, Bollob\'as, Morris, Sahasrabudhe and
Tiba (2020) found signs for which $|P|$ stays between two constant multiples of $\sqrt n$ on the whole circle, settling a question of Littlewood.
\end{solution}

\begin{solution}{pb:skconc}
\idea At high temperature the partition function behaves like its mean: the second moment is at most a constant times the square of the first, so $Z_n\ge\frac12\mathbb EZ_n$
with probability bounded below. Gaussian concentration says $\log Z_n$ hardly moves, so it cannot be below $\log\mathbb EZ_n-O(\sqrt n)$ even with small probability. The
second moment blows up exactly at $\beta=1$.

\writeup (a) $F_n$ is smooth in $g$, with $\frac{\partial F_n}{\partial g_{ij}}=\frac\beta{\sqrt n}\langle\sigma_i\sigma_j\rangle$, where $\langle\cdot\rangle$ is the average under the Gibbs
weights $\frac{e^{H(\sigma)}}{Z_n}$. Since $|\langle\sigma_i\sigma_j\rangle|\le1$, $|\nabla F_n|^2\le\frac{\beta^2}n\binom n2=\frac{\beta^2(n-1)}2$, so $F_n$ is $L$-Lipschitz with
$L=\beta\sqrt{\frac{n-1}2}$ (mean value theorem). Gaussian concentration (Tome~10, Theorem~8.6) gives $P(|F_n-\mathbb EF_n|\ge t)\le2e^{-t^2/2L^2}=2e^{-t^2/\beta^2(n-1)}$.

(b) $H(\sigma)$ is centred Gaussian with variance $\frac{\beta^2}n\binom n2=\frac{\beta^2(n-1)}2$, so $\mathbb Ee^{H(\sigma)}=e^{\beta^2(n-1)/4}$ and $\mathbb EZ_n=2^ne^{\beta^2(n-1)/4}$.
By Jensen, $\mathbb EF_n\le\log\mathbb EZ_n$.

(c) For two configurations $\sigma,\tau$, let $\rho_i=\sigma_i\tau_i$ and $R=\sum_i\rho_i$. $H(\sigma)+H(\tau)$ is centred Gaussian with variance
\[
\frac{\beta^2}n\sum_{i<j}(\sigma_i\sigma_j+\tau_i\tau_j)^2=\frac{\beta^2}n\sum_{i<j}(2+2\rho_i\rho_j)=\frac{\beta^2}n\bigl(n(n-1)+R^2-n\bigr)=\beta^2\Bigl(n-2+\frac{R^2}n\Bigr),
\]
using $2\sum_{i<j}\rho_i\rho_j=R^2-n$. So $\mathbb EZ_n^2=\sum_{\sigma,\tau}e^{\beta^2(n-2)/2+\beta^2R^2/2n}$, and dividing by $(\mathbb EZ_n)^2=4^ne^{\beta^2(n-1)/2}$,
\[
\frac{\mathbb EZ_n^2}{(\mathbb EZ_n)^2}=e^{-\beta^2/2}\,\mathbb E_\rho e^{\beta^2R^2/2n},
\]
where now $\rho$ is a vector of independent fair signs (for uniform independent $\sigma,\tau$, $\rho$ is uniform). Linearise with a Gaussian $g$: $e^{sR^2/n}=\mathbb E_ge^{\sqrt{2s/n}\,gR}$,
and $\mathbb E_\rho e^{bR}=(\cosh b)^n\le e^{nb^2/2}$. So $\mathbb E_\rho e^{sR^2/n}\le\mathbb E_ge^{sg^2}=(1-2s)^{-1/2}$ for $s<\frac12$. With $s=\frac{\beta^2}2$ this is $(1-\beta^2)^{-1/2}$.

(d) By Paley--Zygmund (Tome~10, Lemma~1.13) with $\theta=\frac12$, $P(Z_n>\frac12\mathbb EZ_n)\ge\frac14e^{\beta^2/2}\sqrt{1-\beta^2}=c_\beta>0$. So with probability at least $c_\beta$,
$F_n\ge\log\mathbb EZ_n-\log2$. If $\mathbb EF_n<\log\mathbb EZ_n-\log2-t$ with $t=\beta\sqrt{(n-1)\log(2/c_\beta)}$, then by (a) this event has probability at most
$P(F_n-\mathbb EF_n>t)\le e^{-t^2/\beta^2(n-1)}=\frac{c_\beta}2$, a contradiction. So
\[
n\log2+\frac{\beta^2(n-1)}4-\log2-\beta\sqrt{(n-1)\log(2/c_\beta)}\le\mathbb EF_n\le n\log2+\frac{\beta^2(n-1)}4,
\]
and $\frac{\mathbb EF_n}n\to\log2+\frac{\beta^2}4$. By (a), $P(|F_n-\mathbb EF_n|\ge\varepsilon n)\le2e^{-\varepsilon^2n/\beta^2}$, which is summable, so by Borel--Cantelli
$\frac{F_n}n\to\log2+\frac{\beta^2}4$ almost surely.

(e) $F_n\le n\log2+\max_\sigma H(\sigma)$. The $2^n$ Gaussians $H(\sigma)$ have variance $\frac{\beta^2(n-1)}2$, so by Tome~10, Exercise~8.4 (scaled),
$\mathbb E\max_\sigma H(\sigma)\le\beta\sqrt{\frac{n-1}2}\sqrt{2n\log2}\le\beta n\sqrt{\log2}$. So $\frac{\mathbb EF_n}n\le\log2+\beta\sqrt{\log2}$, which is less than $\log2+\frac{\beta^2}4$ once
$\beta>4\sqrt{\log2}\approx3.33$. The annealed formula of (d) must fail somewhere in $[1,3.33]$. The notebook enumerates all $2^{14}$ configurations: for $n=14$ it finds
$\frac{F_n}n\approx0.752$ at $\beta=0.5$ (annealed value $0.751$), but $1.38$ at $\beta=2$ (annealed value $1.62$).

\tools{Tome~10, Theorem~8.6, Lemma~1.13 and Exercise~8.4; Borel--Cantelli.}
\leads Aizenman, Lebowitz and Ruelle proved in 1987 that the annealed formula holds exactly for $\beta<1$ and fails for $\beta>1$, and found Gaussian fluctuations of $F_n-\mathbb EF_n$ at
high temperature. For $\beta>1$ the limit is the Parisi formula, a variational problem over probability measures on $[0,1]$: Guerra's interpolation (2003) proved the upper bound,
Talagrand (2006) the matching lower bound, and Panchenko (2013) proved the ultrametric structure of the Gibbs measure that Parisi's solution predicted.
\end{solution}

\section*{Map of the chapter}

\begin{center}
\small
\adjustbox{max width=\linewidth}{\begin{tabular}{lll}
Problem & Topic & Tomes used\\\hline
\ref{pb:fixedtail} & tails of the number of fixed points & 3, 4\\
\ref{pb:maxexp} & large deviations of the running maximum & 3, 10\\
\ref{pb:positivewalk} & positive walks: free exponentially, not polynomially & 3\\
\ref{pb:lowentropy} & Sanov for low-entropy words; compression & 3, 10\\
\ref{pb:moderate} & moderate deviations by tilting & 3, 10\\
\ref{pb:windows} & the Erdős--Rényi law of rich windows & 3\\
\ref{pb:lundberg} & Lundberg's ruin exponent, exact & 3, 10\\
\ref{pb:rem} & Derrida's random energy model & 3, 10\\
\ref{pb:markovld} & large deviations for Markov chains; Perron root & 3, 10, 2\\
\ref{pb:discrepancy} & discrepancy: random colourings and Hadamard sets & 3, 2\\
\ref{pb:blowup} & the blowing-up lemma on the cube & 10\\
\ref{pb:withoutrepl} & sampling without replacement is more concentrated & 3\\
\ref{pb:bernstein} & Bernstein's inequality & 3\\
\ref{pb:fpp} & first passage percolation: speed and variance & 10, P1\\
\ref{pb:ksat} & random $k$-SAT: first moment and concentration & 3, 10\\
\ref{pb:janson} & Harris and Janson; triangle-free random graphs & 3, P1\\
\ref{pb:lambdamax} & the top eigenvalue of a GOE matrix & 10, 2\\
\ref{pb:sphere} & Lévy's concentration on the sphere & 10\\
\ref{pb:gumbelmax} & the Gumbel law of a Gaussian maximum & 10\\
\ref{pb:achlioptas} & Johnson--Lindenstrauss with sign matrices & 10\\
\ref{pb:salemzygmund} & random trigonometric polynomials & 4, 10\\
\ref{pb:skconc} & the SK model at high temperature & 10\\
\end{tabular}}
\end{center}

\section*{Sources and further reading}

\begin{itemize}
\item The Nobel Prize in Physics 2021, press release and citation, nobelprize.org; The Abel Prize 2024 citation for Michel Talagrand, abelprize.no. G.~Parisi, Infinite number of
order parameters for spin-glasses, \emph{Physical Review Letters} 43 (1979), 1754--1756; F.~Guerra, Broken replica symmetry bounds in the mean field spin glass model,
\emph{Communications in Mathematical Physics} 233 (2003), 1--12; M.~Talagrand, The Parisi formula, \emph{Annals of Mathematics} 163 (2006), 221--263; D.~Panchenko, The Parisi
ultrametricity conjecture, \emph{Annals of Mathematics} 177 (2013), 383--393.
\item B.~Derrida, Random-energy model: limit of a family of disordered models, \emph{Physical Review Letters} 45 (1980), 79--82; Random-energy model: an exactly solvable model of
disordered systems, \emph{Physical Review B} 24 (1981), 2613--2626.
\item M.~Aizenman, J.~L.~Lebowitz and D.~Ruelle, Some rigorous results on the Sherrington--Kirkpatrick spin glass model, \emph{Communications in Mathematical Physics} 112 (1987),
3--20; M.~Talagrand, \emph{Mean Field Models for Spin Glasses}, 2 vols., Springer, 2011.
\item A.~Dembo and O.~Zeitouni, \emph{Large Deviations Techniques and Applications}, 2nd ed., Springer, 1998 (Markov chains, moderate deviations, G\"artner--Ellis).
\item P.~Erdős and A.~Rényi, On a new law of large numbers, \emph{Journal d'Analyse Mathématique} 23 (1970), 103--111; P.~Deheuvels, L.~Devroye and J.~Lynch, Exact convergence rate in
the limit theorems of Erdős--Rényi and Shepp, \emph{Annals of Probability} 14 (1986), 209--223.
\item F.~Lundberg, \emph{Approximerad framställning af sannolikhetsfunktionen. Återförsäkring af kollektivrisker}, thesis, Uppsala, 1903; H.~Cramér, On the mathematical theory of
risk, \emph{Skandia Jubilee Volume}, Stockholm, 1930; S.~Asmussen and H.~Albrecher, \emph{Ruin Probabilities}, 2nd ed., World Scientific, 2010.
\item J.~Spencer, Six standard deviations suffice, \emph{Transactions of the AMS} 289 (1985), 679--706; N.~Bansal, Constructive algorithms for discrepancy minimization,
\emph{FOCS} 2010, 3--10; S.~Lovett and R.~Meka, Constructive discrepancy minimization by walking on the edges, \emph{SIAM Journal on Computing} 44 (2015), 1573--1582.
\item R.~Ahlswede, P.~Gács and J.~Körner, Bounds on conditional probabilities with applications in multi-user communication, \emph{Zeitschrift für Wahrscheinlichkeitstheorie} 34
(1976), 157--177; L.~H.~Harper, Optimal numberings and isoperimetric problems on graphs, \emph{Journal of Combinatorial Theory} 1 (1966), 385--393.
\item W.~Hoeffding, Probability inequalities for sums of bounded random variables, \emph{Journal of the American Statistical Association} 58 (1963), 13--30; R.~J.~Serfling,
Probability inequalities for the sum in sampling without replacement, \emph{Annals of Statistics} 2 (1974), 39--48.
\item G.~Bennett, Probability inequalities for the sum of independent random variables, \emph{Journal of the American Statistical Association} 57 (1962), 33--45; D.~A.~Freedman,
On tail probabilities for martingales, \emph{Annals of Probability} 3 (1975), 100--118; S.~Boucheron, G.~Lugosi and P.~Massart, \emph{Concentration Inequalities}, Oxford University
Press, 2013.
\item J.~M.~Hammersley and D.~J.~A.~Welsh, First-passage percolation, subadditive processes, stochastic networks, and generalized renewal theory, in \emph{Bernoulli, Bayes, Laplace
Anniversary Volume}, Springer, 1965, 61--110; H.~Kesten, On the speed of convergence in first-passage percolation, \emph{Annals of Applied Probability} 3 (1993), 296--338;
I.~Benjamini, G.~Kalai and O.~Schramm, First passage percolation has sublinear distance variance, \emph{Annals of Probability} 31 (2003), 1970--1978; K.~Johansson, Shape
fluctuations and random matrices, \emph{Communications in Mathematical Physics} 209 (2000), 437--476.
\item E.~Friedgut, Sharp thresholds of graph properties, and the $k$-sat problem, \emph{Journal of the AMS} 12 (1999), 1017--1054; D.~Achlioptas and Y.~Peres, The threshold for random
$k$-SAT is $2^k\log2-O(k)$, \emph{Journal of the AMS} 17 (2004), 947--973; M.~Mézard, G.~Parisi and R.~Zecchina, Analytic and algorithmic solution of random satisfiability problems,
\emph{Science} 297 (2002), 812--815; J.~Ding, A.~Sly and N.~Sun, Proof of the satisfiability conjecture for large $k$, \emph{Annals of Mathematics} 196 (2022), 1--388.
\item T.~E.~Harris, A lower bound for the critical probability in a certain percolation process, \emph{Proceedings of the Cambridge Philosophical Society} 56 (1960), 13--20;
S.~Janson, T.~Łuczak and A.~Ruciński, An exponential bound for the probability of nonexistence of a specified subgraph in a random graph, in \emph{Random Graphs '87}, Wiley, 1990,
73--87; R.~Boppana and J.~Spencer, A useful elementary correlation inequality, \emph{Journal of Combinatorial Theory, Series A} 50 (1989), 305--307; N.~Alon and J.~Spencer,
\emph{The Probabilistic Method}, 4th ed., Wiley, 2016, Chapter~8.
\item C.~A.~Tracy and H.~Widom, On orthogonal and symplectic matrix ensembles, \emph{Communications in Mathematical Physics} 177 (1996), 727--754; S.~Chatterjee,
\emph{Superconcentration and Related Topics}, Springer, 2014.
\item P.~Lévy, \emph{Leçons d'analyse fonctionnelle}, Gauthier-Villars, 1922; M.~Ledoux, \emph{The Concentration of Measure Phenomenon}, AMS, 2001.
\item R.~A.~Fisher and L.~H.~C.~Tippett, Limiting forms of the frequency distribution of the largest or smallest member of a sample, \emph{Proceedings of the Cambridge Philosophical
Society} 24 (1928), 180--190; B.~Gnedenko, Sur la distribution limite du terme maximum d'une série aléatoire, \emph{Annals of Mathematics} 44 (1943), 423--453.
\item D.~Achlioptas, Database-friendly random projections: Johnson--Lindenstrauss with binary coins, \emph{Journal of Computer and System Sciences} 66 (2003), 671--687; D.~M.~Kane
and J.~Nelson, Sparser Johnson--Lindenstrauss transforms, \emph{Journal of the ACM} 61 (2014), 4:1--4:23; K.~G.~Larsen and J.~Nelson, Optimality of the Johnson--Lindenstrauss lemma,
\emph{FOCS} 2017, 633--638.
\item R.~Salem and A.~Zygmund, Some properties of trigonometric series whose terms have random signs, \emph{Acta Mathematica} 91 (1954), 245--301; G.~Halász, On a result of Salem
and Zygmund concerning random polynomials, \emph{Studia Scientiarum Mathematicarum Hungarica} 8 (1973), 369--377; P.~Balister, B.~Bollobás, R.~Morris, J.~Sahasrabudhe and
M.~Tiba, Flat Littlewood polynomials exist, \emph{Annals of Mathematics} 192 (2020), 977--1004.
\end{itemize}
